% prime_ext.tex -- prime-field expansion: three chapters that follow
% Chapter~\ref{sec:prime} (sections/prog_prime.tex).
% Labels created: ch:primegeneric, ch:primealgebraic, ch:primetargets and
% the sec:pg-*, sec:pa-*, sec:pt-*, tab:pg-*, tab:pa-*, tab:pt-* labels below.

\chapter{Prime fields: the generic track}\label{ch:primegeneric}

Chapter~\ref{sec:prime} gave the prime-field programme as it stood in
mid-October: the known-answer ladders, Table~B, the Kummer tower, the
structured P-256 bases, the isogeny-class walks and the lattice
reformulations. This chapter and the next two record what happened after
the owner's decisions of 29 October split prime-field index calculus into
two counted tracks (\art{research/prime_ic_program_20261029/DECISIONS.md},
decision R1; the split itself is stated in Chapter~\ref{ch:distance},
Table~\ref{tab:distance-prime}). The \emph{generic} track holds every
oracle built from tables, pair tables, root finding, subtraction or
collision search. Its improvements count and ratchet its own baseline, but
they carry the label ``bounded by $\sqrt r$; cannot beat rho''. The
\emph{algebraic} track, the subject of Chapter~\ref{ch:primealgebraic},
holds relations from a solver that receives the factor-base constraint
inside its system.

The chapter has one result that matters for the distance to rho and one
correction. The residual collectors of \S\ref{sec:pg-residual} are the
prime-field arms closest to rho that the programme measured: $1.54$--$1.64\times$
matched rho at $w=0$ on 36--37-bit curves with a validator report, and
$1.3$--$1.7\times$ on 36--42-bit curves as the producer reports them, with no
growth of that distance over the measured sizes. The Curve25519 walk tables
of \S\ref{sec:c25519} are at about the same distance ($1.6$--$1.9\times$,
unverified). They are generic collision searches with rho's
exponent, not index calculus with a better one. Their large gains at
$w>0$ were a pricing artefact, withdrawn when the owner chose option~B on
1 November (\S\ref{sec:pg-bytes}).

\paragraph{Standing of the numbers.} Under the rule of 4 November every
producer is assumed to be possibly sabotaging the work, and a result is
cited by its proof-carrying claim and standing or labelled an unverified
assertion. None of the generic-track studies in this chapter has a
\texttt{pcc} claim record. Four evidence directories have a
\texttt{VALIDATION.json} with a replay certificate and a validator
report: Table~B (Chapter~\ref{sec:prime}), the 20-bit $N=45$ baseline
(Chapter~\ref{ch:distance}), the first $m=4$ cells of
\S\ref{sec:pg-gate}, and the packed membership filter of
\S\ref{sec:pg-engineering}. The residual collector study of 1 November has
a replay certificate, a validator report and a \texttt{VALIDATION.json},
but is not an admitted row (\S\ref{sec:pg-residual}). Every other figure
below is a producer-reported value from a committed record, re-read from
the source for this chapter but not independently replayed, and is so
labelled where it is used. Costs are in $S=(\text{counted GAE}+w\cdot
\text{unpriced native operations})/\sqrt r$ at the band weights
$w=0,0.1,1$ of \S\ref{sec:band}, and two arms are ranked only where their
order holds across the band.

\paragraph{The programme of 29 October.} The decisions file that created
the tracks also set the programme's execution plan
(\art{research/prime_ic_program_20261029/DECISIONS.md}). Its entries are,
in brief:
\begin{itemize}
\item \texttt{D-EKS}: compute for the published algebraic reference arm,
  a derived estimate of about 49 CPU-hours (range 15--150) on EKS Spot,
  approved;
\item \texttt{D2}: the Table~B incumbent may be ranked against that
  reference, which amends R1 so that arms of the two tracks may be ranked on
  the same instances, in the same unit and across the band, while each track
  ratchets only its own baseline;
\item \texttt{D3}: a pinned, vendored msolve counter patch is allowed, its
  hash bound into the candidate identity;
\item \texttt{D-PREFETCH}: the hash-table layout measurement of
  \S\ref{sec:pg-engineering}, approved;
\item \texttt{D-FD7} and \texttt{D-FD8}: two amendments of the prime
  fall-degree study, the two-cap bracket for $t\ge12$ with 56/112\,GiB jobs
  and the $m=4$ extension, adopted;
\item \texttt{D-MAC}: the fall-degree study's raw runs, which lived only in
  one worktree on the owner's laptop, are to be recovered rather than
  declared lost; until then its $t\le10$ reproduction stays at median level;
\item \texttt{D-ROUTE}: a command route for non-ecbench jobs on the lab,
  approved for an administrator to implement;
\item \texttt{D-REPLAY}: a second-host replay of the 20-bit holdout session,
  then a claims record, approved.
\end{itemize}
The same file records why root finding is generic: at $m=2$,
root-find-and-subtract fixes $x_1=c$ for each $c\in V$, solves the remaining
quadratic $\Sem3(c,X,x_R)$ by its roots $x(R\mp P_c)$ and tests each root
against a hash index of $V$. It returns the solution set of a Gr\"obner solve
of $\langle\Sem3,F_V(x_1),F_V(x_2)\rangle$, but its cost never depends on
$F_V$, and it runs identically on a random $V$.

Table~\ref{tab:pg-summary} lists the prime-field arms of this chapter and
of Chapter~\ref{sec:prime} that have a distance to matched rho, from the
closest to the farthest.

\begin{table}[t]
\centering\small
\begin{tabularx}{\linewidth}{@{}Xlll@{}}
\toprule
Arm & Bits & IC/rho ($w=0$) & Standing \\
\midrule
Residual collector \texttt{res1eq} (\S\ref{sec:pg-residual}) & 36--42 & $1.32$--$1.72$ & validator report at 36--37 bits ($1.54$--$1.64$); producer at 36--42 bits \\
Residual collector \texttt{res1dp} & 36--42 & $1.92$--$2.27$ & as above \\
Residual collector \texttt{res2} & 36 & $2.6$ & improvement-log row, no validator \\
Residual collector \texttt{res4} & 20.6--32 & $3.78$--$4.94$ & producer \\
$k=5$ sum table (\S\ref{sec:pg-ksum}) & 20--48 & $5.4$--$11.0$ & unqualified \\
Table~B pair table (Chapter~\ref{sec:prime}) & 10--24 & $5.47$--$13.5$ & validated \\
20-bit pair table, $N=45$ (Chapter~\ref{ch:distance}) & 20 & $7.75$ & validated \\
Cofactor-cleared base, $h=32$ (\S\ref{sec:pg-cofactor}) & census & $9.80$ & \texttt{invalid\_controls} \\
$k=4$ sum table (Chapter~\ref{ch:distance}) & 20--40 & $5.3$--$13.1$ & no validation \\
Kummer tower, P-224 shape (\S\ref{sec:pa-pkm}) & 20--28 & $18$--$78$ & producer \\
Kummer tower (Chapter~\ref{sec:prime}) & 20--30 & $29$--$159$ & producer \\
Packed pair table with filter (\S\ref{sec:pg-engineering}) & 32 & $31.47$ & validated with caveats \\
\bottomrule
\end{tabularx}
\caption[]{Prime-field index calculus against matched negation rho,
greater than 1 meaning rho is cheaper, as each source states it. Ranges are
over sizes or curves; each row's band behaviour is in its section. The
residual rows are at rho's exponent; every other row's gap grows with
size where more than one size was measured.}
\label{tab:pg-summary}
\end{table}

\section{What a generic oracle can buy}\label{sec:pg-gate}

\paragraph{The arity gate.} Let $s=\log_2 p/m$, so that a factor base of
the size index calculus needs has $\abs{V}\asymp 2^s$. Suppose a
target-independent setup costs $2^{as}$ and one natural decomposition query
costs $2^{es}$. Collecting $\Theta(2^s)$ relations costs
$2^{as}+2^{(e+1)s}$ up to polynomial factors, so beating rho's $p^{1/2}$
needs both
\[
  a<\frac m2 \qquad\text{and}\qquad e<\frac m2-1 .
\]
At $m=2$ the second condition is $e<0$, impossible for a nontrivial query,
so exact pair enumeration is optimal at that arity up to polynomial
factors; the first arity that can discriminate is $m=3$, with $a<1.5$ and
$e<0.5$. The derivation is the arity and solver-model correction of 31
October and 1 November (\texttt{DEC-20261031-b51282},
\texttt{DEC-20261101-26f1fa};
\art{research/prime_ic_lead_review_20261031/PDP_ARITY_CORRECTION.md}),
which found that every accepted theorem and completed measurement of the
prime factor-base degree study was a two-summand result. The same pair of
gates is the ``complete-cost'' row of the six-obligation chain of
\S\ref{sec:pg-obligations}.

\paragraph{A wrong gate and its replacement.} An audit of 4 November of the
programme memo \emph{Alternatives to Semaev summation polynomials for
prime-field ECDLP index calculus}
(\art{research/semaev_alternatives_audit_20261104/}) found that the memo's
oracle gate, ``$e<1$ at $m=4$'', is not implied by its own cost model. A
generic 2+2 meet-in-the-middle oracle passes it, with measured $s=1.93$ on
$T=B^s$ and $e=0.4$ at $B=p^{1/5}$, while costing $p^{2/3}$, more than rho.
The audit's replacement is $s<m-3$, that is $e<(m-3)/4$. Seven Lean~4
theorems in \art{research/semaev_alternatives_audit_20261104/lean/}
(\texttt{Routes.lean}, \texttt{Roots.lean}) cover arities 2 and 3, the
$m=4,5$ gate, the meet-in-the-middle route and a two-roots bound; the
audit's own ledger counts only these theorems as ``ruled out''. Its
measured numbers are a single local run with no replay certificate, and
its README says to treat them as unverified until re-run under the
validation gate. The memo's ``Idea~10'' measured one squarefree collision
for all $B\le2^{16}$; Ideas~2 and~3 are $m=2$ objects and cannot change the
exponent.

\paragraph{The $m=4$ planted solve.} The gate correction has a measured
counterpart. A direct 2+2 meet-in-the-middle at $m=4$ over a random base of
size $B=\mathrm{round}(p^{1/5})$ (\art{research/ic_m4_mitm_scaling_20261031/};
stage diagnostic, producer-reported) recovered every planted four-term
decomposition with a fitted exponent of $0.412$ in $\log_2 p$ ($R^2=0.9986$,
eight sizes to 52 bits), $0.404$ on a three-rung ladder to 60-bit primes
and $0.396$--$0.450$ on an arbitrary-precision ladder to 70 bits. The
planted solve is not a relation collector. A uniform target decomposes
with probability about $B^4/p$, so one relation costs about $p/B^2=p^{3/5}$:
the 24--40-bit fit of the uniform arm is $0.469$, but the local slope from
40 to 44 bits is $\log_2(1.08\cdot10^9/2.09\cdot10^7)/4=1.41$, and the arm
was stopped there. The README's own reading is that the $m=4$ solve ``is not
an end-to-end sub-1/2 relation collector'', which is the audit's point:
an oracle exponent below $1/2$ in the planted regime says nothing about
the cost of a relation.

\paragraph{The first $m=4$ cost cells.} The blind-spot programme of 18
October priced arity $m=4$ end to end on 16--28-bit primes
(\art{benchmarks/20261018-blindspots-prime-m4-cells/}; validated, with a
negation-map rho added to the same binary and the index-calculus rows
reproduced exactly in it). Full index calculus costs $8.2$--$29.5\times$ the
negation-map rho at $w=0.1$ and $6.5$--$35.6\times$ across the band, dearer
than both rhos at every cell and weight, 16 of 16 solves per cell; the
arm's own F4 grows as $\abs{\Fb}^{11.03}$. The verdict is \cls{regresses}
at the primitive level.

\paragraph{Exponents on the bounds frontier.} The sealed bounds frontier
(\art{docs/bounds/FRONTIER.md}, toy tier, 17.7--24.0 bits, 96 runs) gives
the folded pair-table \texttt{ic.pipeline} a fitted exponent of $0.841$
$[0.646,0.953]$ and the subtraction oracle $1.005$ $[0.880,1.077]$, against
$0.504$ $[0.316,0.585]$ for interleaved baby-step giant-step. These are
fit-version-2 bounds with wide intervals at four sizes; the campaign
exponent $0.620$ $[0.557,0.663]$ of Table~\ref{tab:prog_frontier-exponents}
is the better-determined estimate of the same pair table.

\paragraph{Arity sweep.} The crypto-autoresearcher index-calculus sweep of
27 October (\art{research/ic_arity_sweep_20261027/}; 234 rows, every
logarithm checked) ran $m=2,\dots,9$ on generated prime-order curves at
12--18 bits with exhaustive and meet-in-the-middle engines. Its unit is the
tool's own, $\Sem3$ root solves against rho walk additions, which the
tool says are different units. The best exhaustive arity is $m=2$ at
$10.6$, $16.2$, $20.4$ and $60.5\times$ rho at 12, 14, 16 and 18 bits; the
best meet-in-the-middle arity is $m=5$ or $m=3$ at $2.8$, $3.3$, $5.8$ and
$4.3\times$. Odd arities have the lower fitted exponents ($0.63$ at $m=3$,
$0.60$ at $m=9$) because the default table split leaves one more level on
the search side at even $m$. None of this is in the counted ecbench unit,
and none of it is below rho.

\section{$k$-sum tables beyond $k=4$}\label{sec:pg-ksum}

Chapter~\ref{ch:distance} gave the $k=4$ sum table of 28 October: a
negation-folded table of every signed 4-sum of a smaller base, $S$ $83.7\to
11.6$ at 40 bits, $13.1\times$ \texttt{rho-neg} there. Its generic law is
$2^{k-1}\binom{A}{k}+c_k A p/(2^k\binom Ak)$, with optimum exponent
$k/(2k-1)$, falling toward $1/2$ (baby-step giant-step) as $k$ grows and
never below it. Three studies followed (Table~\ref{tab:pg-ksum}). All three
are registered before data, every run verified and replay-audited,
and none has a second-host replay, a \texttt{VALIDATION.json} or a
\texttt{claims/*.json} record; the two later studies call themselves
\emph{unqualified}. The benchmark index lists them
only as ``not admitted to the results table until it has an
independent-host replay and its own \texttt{VALIDATION.json}''.

\begin{table}[t]
\centering\small
\begin{tabular}{@{}lllll@{}}
\toprule
Study & Bits & Ratio to its incumbent ($w=0$) & To \texttt{rho-neg} & Exponent \\
\midrule
$k=5$ vs $k=4$ & 40 & $0.74$ ($0.71$ at $w=1$) & $9.7\times$ & $0.547$ $[0.543,0.554]$ \\
 & 48 & $0.67$ & $11.0\times$ & ($k=4$: $0.565$) \\
\addlinespace
$k=5$ vs mitm incumbent & 20 & $0.701$ $[0.661,0.744]$ & $5.4\times$ & \\
(same instances) & 24 & $0.453$ $[0.433,0.475]$ & $6.3\times$ & \\
 & 28 & $0.331$ $[0.320,0.344]$ & $7.2\times$ & \\
 & 32 & $0.218$ $[0.211,0.226]$ & $7.3\times$ & \\
\addlinespace
$k=5$ vs re-tuned incumbent & 20 & $0.700$ $[0.682,0.718]$ & $6.0\times$ & $0.496$ $[0.487,0.510]$ \\
(16 fresh curves) & 24 & $0.446$ $[0.432,0.461]$ & $5.8\times$ & (incumbent: \\
 & 28 & $0.324$ $[0.318,0.331]$ & $7.3\times$ & $0.635$ $[0.627,0.645]$) \\
 & 32 & $0.217$ $[0.212,0.222]$ & $7.3\times$ & \\
\bottomrule
\end{tabular}
\caption[]{Prime-field $k$-sum tables after 28 October (generic track;
all unqualified, producer-reported). Ratio: paired ratio of total cost,
cluster-bootstrap 95\% intervals. ``To rho'': mean $S$ over matched
negation rho at $w=0$. Exponent: fitted cost exponent in $r$ at $w=0$;
the pooled study's rho-neg fits $0.457$ and \texttt{bsgs-neg} $0.501$ on
the same sessions. Sources: \art{research/prime_ksum_k5_20261028/},
\art{research/prime_ksum_vs_incumbent_20261029/},
\art{research/prime_ksum_pooled_20261029/}.}
\label{tab:pg-ksum}
\end{table}

\paragraph{$k=5$.} A table of every signed 5-sum, with indices packed 12
bits each into one word and the insert order unchanged so that
group-operation counts match, lowers $S$ below the $k=4$ table at 36 and 40
bits at every band weight, the registered success test
(\art{research/prime_ksum_k5_20261028/}; 1{,}920 runs, all eight test
sessions reproduce every replay; $k=4$ reproduces the earlier per-run
counts on 288 of 288 runs). At 40 bits $S$ is $8.51$ against $11.56$. The
fitted exponent over 36--48 bits is $0.547$ $[0.543,0.554]$ against $0.565$
$[0.560,0.571]$, finite-size fits near the generic optima
$5/9\approx0.556$ and $4/7\approx0.571$. The table is still $7.6\times$ baby-step giant-step at 40 bits and
$10.6\times$ at 48, so the method sits above its own generic limit at every
size.

\paragraph{Against the incumbent.} The tracks rule requires a same-instance
ranking before a generic change replaces the incumbent. On the incumbent's
own four holdout curves, 48 targets, method hashes and measurement seed
(\art{research/prime_ksum_vs_incumbent_20261029/}; 3{,}456 runs) the $k=5$
table costs less than the negation-folded pair table of
\art{research/prime_base_oracle_tuning_20261029/} at every size and every
weight, and the registered primary test holds in all three of its cells
(24, 28, 32 bits). The generic baseline on those curves falls by $1.4\times$
at 20 bits and $4.6\times$ at 32. The widening gap is what the two optima
predict, $5/9$ against about $2/3$; one curve per size gives no exponent.
The study compared against the recorded incumbent, whose 32-bit base was
null by construction, so a second study re-tuned the incumbent on separate
train curves by the same rule as the $k$-sum families and tested both on 16
fresh curves (\art{research/prime_ksum_pooled_20261029/}; 13{,}824 runs).
The ranking survived: $k=5$ is below the tuned incumbent on all 16 curves
and in the pooled test at all three weights. Its toy-tier exponent of
$0.496$ is a finite-size fit over 20--32 bits; the generic law does not
allow a limit below $1/2$. Tuned for $w=0$, the incumbent
is 8--13\% cheaper, and $k=5$ is still $2.2$--$4.6\times$ cheaper than it. A train-frozen
rule chose $k=5$ at every size; at 20 bits $k=4$ is cheaper at $w=0$ and
$0.1$ and the two are unresolved at $w=1$.

\paragraph{The precomputation regime.} The $k=4$ table is built before any
target arrives, which invites two accounting questions
(\art{research/prime_precompute_regime_20261028/}; class: control and
accounting, not an improvement). At equal memory against preprocessing rho
(Bernstein and Lange~\cite{bl2012}), a known-log variant
\texttt{kl-ksum4} answers each target with one probe and no matrix. At 40
bits the equal-size preprocessing rho's online $S$ is below
\texttt{kl-ksum4}'s by $142$, $157$ and $229\times$ at $w=0,0.1,1$, and its
precomputation costs $333$--$339\times$ more. Preprocessing rho sits on
$E\cdot T^2\approx N$ ($6.96$ at $w=0$); the known-log table sits about
$N/E\approx1.4\cdot10^5$ above that line, as the Corrigan-Gibbs--Kogan bound
$ST^2\gtrsim N$~\cite{cgk2018} allows. Table~\ref{tab:pg-precompute} compares the registered
predictions with the 40-bit measurement. The amortised crossover between the
two falls at $14.3$k, $13.1$k and $9.5$k targets at 40 bits
($w=0,0.1,1$), against a registered prediction of $15$k. Against
preprocessing rho at equal precomputation there is no positive crossover
at 36 and 40 bits. These are multi-target and preprocessing numbers; no
single-target cost in the study falls below rho's.

\begin{table}[t]
\centering\small
\begin{tabular}{@{}lrr@{}}
\toprule
Quantity (40 bits, $w=0$) & Predicted & Measured \\
\midrule
\texttt{kl-ksum4} set-up $S$ & 5.3 & 5.01 \\
\texttt{kl-ksum4} online $S$ & 0.10 & 0.117 \\
equal-size preprocessing rho, set-up $S$ & 1{,}594 & 1{,}668 \\
equal-size preprocessing rho, online $S$ & 0.0006 & 0.0008 \\
equal-precomputation preprocessing rho, online $S$ & 0.19 & 0.33 \\
\texttt{kl}/equal-size crossover (targets) & 15k & 14.3k \\
\bottomrule
\end{tabular}
\caption[]{Registered predictions against measurement for the known-log
$k=4$ table and preprocessing rho at 40 bits (class: control and
accounting). Producer-reported. Source:
\art{research/prime_precompute_regime_20261028/}.}
\label{tab:pg-precompute}
\end{table}

\paragraph{Known-log targets.} A target-blind known-log two-sum method of
21 October (\art{research/ecbench_known_log_targets_20261021/}) recovered
public targets without receiving their logarithms. On 16-, 20- and 24-bit
fixtures its cold counted cost was below the abscissa pipeline's on both
the primary and holdout panels at every band weight; at 24 bits the
pipeline/known-log ratio is $15.6$--$17.0$ on the primary panel. Negation
rho remained cheaper at every size ($S$ $1.16$ against $2.70$ at 24 bits),
the 12-bit comparison is unresolved inside the band, wall measurements are
L0 and the receipts are same-host replays. Table~\ref{tab:pg-knownlog} gives the
primary panel.

\begin{table}[t]
\centering\small
\begin{tabular}{@{}lrrrl@{}}
\toprule
Curve ($\log_2r$) & Abscissa $S$ & Known-log $S$ & Rho $S$ & Abscissa / known-log, $w=0,0.1,1$ \\
\midrule
\texttt{icv1-fp12-tm53} (11.807) & 27.78 & 25.06 & 2.92 & $1.109$, $1.049$, $0.755$ \\
\texttt{icv1-fp16-t295} (15.697) & 11.10 & 7.57 & 1.76 & $1.467$, $1.447$, $1.348$ \\
\texttt{icv1-fp20-t727} (19.101) & 12.09 & 2.96 & 1.36 & $4.081$, $4.295$, $5.336$ \\
\texttt{icv1-fp24-t1059} (23.974) & 42.23 & 2.70 & 1.16 & $15.65$, $15.96$, $16.98$ \\
\bottomrule
\end{tabular}
\caption[]{Target-blind known-log two-sum recovery against the abscissa
pipeline ($32$ unsigned pairs, folded two-sum meet-in-the-middle) and
negation rho, cold, eight targets per curve and five rounds; $S$ at $w=0$.
Producer-reported, same-host replays, no validation. Source:
\art{research/ecbench_known_log_targets_20261021/}.}
\label{tab:pg-knownlog}
\end{table}

\section{Residual collectors}\label{sec:pg-residual}

A partial decomposition $R=P_u\pm P_v$ with $P_v\in\Fb$ leaves a residual
$x(R\mp P_v)$ that misses $\Fb$. The residual collector
(\texttt{ic\_framework::residual}, \texttt{targets=residual}) keeps every
residual and turns residual collisions into relations; it is the
$\abs{D}=r$ limit of the large-prime variation. The programme ran it on
prime curves from 26 October to 2 November, under the atlas lead
\texttt{LEAD-ICP-residual-partial-relations} and hypothesis
\texttt{H-PFDR-cb9606}. Table~\ref{tab:pg-residual} gives the arms nearest
rho. A residual collision is a birthday collision on $\ZZ/r$ with a base
attached; the collector's cost exponent matches rho's, and on a curve-free
carrier it costs the same as on the curve. That is the content of
\texttt{KeyedInvariance.lean}
(\art{research/ic_planner_exhaustive_20261026/lean/}): an algorithm that
sees elements only through a keyed labelling cannot use the curve.

\begin{table}[t]
\centering\small
\begin{tabularx}{\linewidth}{@{}llrrX@{}}
\toprule
Arm & Curve & $w=0$ & $w=1$ & Source and standing \\
\midrule
\texttt{res4} (base 4) & 20.6--32 bits & $3.78$--$4.94$ & $3.6$--$4.8$$^\ast$ & residual\_lp; producer \\
\texttt{res2} & p36 (36.0 bits) & $2.6$ & $2.6$$^\ast$ & planner\_exhaustive; log row $2.59\times$ \\
\texttt{res2} & p37 (36.25 bits) & $2.7$ & $2.7$$^\ast$ & planner\_exhaustive; producer \\
\texttt{res4f} (fused) & p36 / p37 & $2.39$ / $2.51$ & --- & residual\_fused; producer \\
\addlinespace
\texttt{res1eq} & p36 & $1.54$ & $2.66$ & residual\_dp\_eq; validator report \\
\texttt{res1eq} & p37 & $1.64$ & $2.83$ & residual\_dp\_eq; validator report \\
\texttt{res1dp} & p36 / p37 & $2.13$ / $2.03$ & $2.22$ / $2.07$ & residual\_dp\_eq; validator report \\
\addlinespace
\texttt{res1eq} & p36 / p40 / p42 & $1.32$ / $1.72$ / $1.39$ & $2.27$ / $2.99$ / $2.42$ & residual\_scale; producer \\
\texttt{res1dp} & p36 / p40 / p42 & $2.15$ / $2.27$ / $1.92$ & $2.25$ / $2.28$ / $1.93$ & residual\_scale; producer \\
\texttt{res1eqc} (chained) & p36 & $1.591$ & $2.73$ & res1eq\_chain; producer \\
\bottomrule
\end{tabularx}
\caption[]{Residual collectors: distance to matched negation rho
(cost over rho, greater than 1 means rho is cheaper). $^\ast$\,With the
table-byte counter excluded, as option~B prices the band
(\S\ref{sec:pg-bytes}); the as-published $w=1$ values were $41$--$58\times$
(\texttt{res4}) and $23$--$24\times$ (\texttt{res2}). The \texttt{eq} and
\texttt{dp} rows are nearly unchanged by option~B. Sources under
\art{research/}: \texttt{ic\_prime\_residual\_lp\_20261026},
\texttt{ic\_planner\_exhaustive\_20261026},
\texttt{ic\_residual\_fused\_20261031},
\texttt{ic\_residual\_dp\_eq\_20261101}, \texttt{ic\_residual\_scale\_20261101},
\texttt{ic\_res1eq\_chain\_20261101}; validator report
\art{benchmarks/20261101-ic-residual-dp-eq/validator-report.json}.}
\label{tab:pg-residual}
\end{table}

\paragraph{The first measurement.} On six registered prime curves of
20.6--32 bits, against the \texttt{ic plan} recipe and \texttt{rho.negation}
(\art{research/ic_prime_residual_lp_20261026/}; 24 of 24 runs verified per
arm and curve), the collector with a four-point base was $2.3$--$8.4\times$
cheaper than the plan in counted $S$ (\texttt{res4}/plan $0.119$--$0.426$ at
$w=0$, every interval below 1). As published, the same arm cost
$1.33$--$3.89\times$ the plan at $w=1$, so the order reversed inside the
band at every size and the two were not ranked. The least-squares exponents
over six sizes are $0.470$ for \texttt{res4} and $0.480$ for rho, against
$0.634$ for the plan; at any fixed base the collector grows as about
$r^{1/2}$. On the curve-free carrier ($\ZZ/r$ with a random-bijection key,
same seeds) all nine intervals contain 1. The collector buys its group-operation
savings with a table probe and an insert on every step, so the native work
it adds is real work; what option~B later removed from the $w>0$ figures
was a byte count, not those probes (\S\ref{sec:pg-bytes}). The $k$-sum
study derived that large primes do not help the \emph{two-summand pair
table}: with two summands a large-prime table of
$M$ entries has the same law $M+cp/\sqrt M$ (\art{research/prime_ksum_table_20261028/}).
The residual collector escapes that law because it gives up the pair table
and becomes a collision search.

\paragraph{The planner re-test.} The study extrapolated a crossing of the
plan near 35 bits. The planner re-test of the same day
(\art{research/ic_planner_exhaustive_20261026/}; 14 new sessions, 3{,}328
records, all verified; 2{,}428 band verdicts restated in integers and
checked by Lean's \texttt{decide}) measured it. At $36.0$ and $36.25$ bits
\texttt{res2} and \texttt{res4} are cheaper than the plan at every band
weight with intervals below 1 (\texttt{res4}/plan $0.888$ $[0.817,0.965]$
and $0.865$ $[0.802,0.933]$ at $w=1$, the risky registered prediction R3).
Over all ten sizes the $w=1$ ratio falls steadily, from $4.92$ at 15.7 bits
to $0.865$ at 36.25. The curve-free null held (15 of 16 intervals contain 1)
and the fitted slope over ten sizes is $0.460$ for \texttt{res4} against
$0.619$ for the plan. At 36 bits the collector moved the framework's
cheapest one-target arm from $41$--$47\times$ rho in counted $S$ to
$2.6$--$2.7\times$. The leaderboard improvement log carries this as
``Residual collector res2 vs rho.negation, fp36: $2.59\times$''
(\art{docs/ic/LEADERBOARD.md}), the $w=0$ value; the study's own decision
notes that no validator has run on it, so it is an unvalidated row.

\paragraph{Fused steps, distinguished points and the equivariant walk.}
The fused residual step of 31 October
(\art{research/ic_residual_fused_20261031/}; four sessions of 512 records,
all verified, A/A $1.0000$) cut group additions per step from $2.005$ to
$1.004$ and gave fused/baseline $0.56$--$0.92$ at $w=0$ on all 12
(curve, base) cells. The fused walk needs $1.1$--$1.3\times$ more steps per
harvested relation, so priced native work ate the saving: the order is
mixed at $w=0.1$ and reversed at $w=1$, and the preregistered rule C1 was
not met. Distinguished points and an equivariant walk followed on 1
November (\art{research/ic_residual_dp_eq_20261101/}; two sessions of 704
records, all verified, A/A $1.0000$; it supersedes the $w>0$ figures of
\art{research/ic_residual_dp_20261101/}, whose per-step distinguished-point
test was unpriced). The best arms are \texttt{res1eq} at $1.54$, $1.75$ and
$2.66\times$ matched rho on p36 ($w=0,0.1,1$) and $1.64$, $1.86$, $2.83$ on
p37, and \texttt{res1dp} at $2.13$, $2.14$, $2.22$ and $2.03$, $2.04$,
$2.07$. Between \texttt{eq} and \texttt{dp} the order is unresolved inside
the band: \texttt{eq} is lower at $w=0$, \texttt{dp} at $w=1$.

\paragraph{Scale.} On p36, p40 and p42 (\art{research/ic_residual_scale_20261101/};
three sessions of 288 records, 863 verified, one baseline run exhausted at
$10^8$ trials) the distance of \texttt{res1eq} to matched rho is $1.32$,
$1.72$ and $1.39$ at $w=0$ and $2.27$, $2.99$, $2.42$ at $w=1$; of
\texttt{res1dp} $2.15$, $2.27$, $1.92$ and $2.25$, $2.28$, $1.93$. There is
no growth of the distance with $r$ over this range, as a method with rho's
exponent predicts. Chained restarts, which restart each walk at the
previous start plus one random jump instead of two scalar multiplications,
were null: $0.994$ $[0.953,1.037]$ at $w=0$
(\art{research/ic_res1eq_chain_20261101/}).

\paragraph{Batches and many targets.} A residual batch of $L$ targets
against pooled multi-target rho at p36
(\art{research/ic_residual_batch_ecbench_20261102/}; 320 of 320 verified,
local L0, L2 replay owed) holds the frozen parity rule at $L=64$, $1.0077$
$[0.9811,1.0350]$, and fails it at $L=16$, $1.0465$ $[0.9886,1.1078]$. Per
target both batch methods are $3.3$--$3.4\times$ ($L=16$) and
$6.6$--$6.7\times$ ($L=64$) cheaper than single-target rho; the matched
reference for a batch is multi-target rho, and against it the collector is
at parity, not ahead. An exploratory shared-table collection on a 20-bit
curve (\art{research/ic_largeprime_multitarget_20261102/}; not
preregistered) gave a per-target cost at $T=64$ targets $4.4\times$ below
$T=1$. That is multi-target rho reproduced. The same note sets aside
triple and quadruple large primes as a generic $k$-sum search already
measured in \S\ref{sec:pg-ksum}.

\subsection{Correction: the byte counter}\label{sec:pg-bytes}

The independent validator of the 1 November study
(\art{benchmarks/20261101-ic-residual-dp-eq/validator-report.json},
finding F1) found that the band as built priced
\texttt{residual\_table\_bytes\_uncharged}, a byte count of the residual
table, as native operations at $w>0$. The audit of the same day
(\art{research/ic_byte_counter_pricing_20261101/}) confirmed it: the
counter is emitted in \texttt{ic\_framework/residual.rs}, every
\texttt{*\_uncharged} counter is appended to the unpriced list, and
\texttt{priced\_gae(w)} charges each listed unit $w$ GAE, so each byte cost
$w$ GAE. The analogous \texttt{native\_table\_bytes\_uncharged} is classed
as shape, not work. The counter was 83--99\% of the residual arms' priced
$S$ at $w=1$ and 40--98\% at $w=0.1$, and it is the whole reason their $w>0$
distances to rho read $17$--$58\times$ instead of $2.6$--$5\times$. On the
first residual study, \texttt{res4}/rho at $w=1$ falls from $41.2$--$57.6$
as published to $3.6$--$4.8$ with the counter excluded; on the planner
re-test, \texttt{res2}/rho at p36 falls from $23.3$ to $2.6$.

The owner chose option~B on 1 November: the band prices operations, not
bytes. \texttt{band.rs} now lists the byte counters as
\texttt{NON\_OPERATION\_COUNTERS} and does not price them; memory is
reported on its own axis. No session record was edited; the four affected
studies carry a dated note and the re-run comparisons under
\texttt{analysis/exbytes/}. What changed:
\begin{itemize}
\item \emph{The large $w>0$ gains are withdrawn.} They were the byte price.
  On the planner re-test, 36 of 136 weighted comparisons go from
  ``unresolved inside the band'' to resolved, 12 stay unresolved and 88 keep
  their direction; the $w>0$ distances are flat in $w$.
\item \emph{The rule C1 of the scale study is met in 2 of 8 cells, not 6 of
  8.} The equivariant arm fails all four cells at $w=1$, where its interval
  reaches or crosses 1 (p40 \texttt{res1eq} $1.181$ $[0.859,1.625]$).
\item \emph{An unreconciled tally.} The 1 November errata give the 12-cell
  C1 tally as 9 of 12; the study's own \texttt{analysis/summary.json} gives
  8 of 12 as built. Nobody has reconciled the ninth cell.
\item \emph{What stands.} $w=0$ results, the leaderboard's $2.59\times$,
  every verification, audit and A/A result, and the absence of growth of the
  distance with $r$.
\end{itemize}
A second byte counter, \texttt{factor\_base\_vec\_bytes\_floor\_uncharged},
is priced the same way in the Koblitz sessions, at $0.4$--$12\%$ of $S$ at
$w=1$; no Koblitz verdict changes direction. The 1 November evidence
directory has a replay certificate, the validator report and a
\texttt{VALIDATION.json}, but is not an admitted row: its raw records live
in the artifact store, which the CI gate cannot read
(\art{benchmarks/20261101-ic-residual-dp-eq/README.md}).

\paragraph{What the collectors are.} The residual collectors are the
nearest the programme came to rho on a prime field: $1.3$--$1.7\times$ at
$w=0$ and $2.2$--$3.0\times$ at $w=1$ for \texttt{res1eq}, and about
$2\times$ at every weight for \texttt{res1dp}. They are measured constant
improvements of the generic track over the pair-table incumbent, at rho's
exponent. By the curve-free null and \texttt{KeyedInvariance.lean} they
are birthday searches on $\ZZ/r$, and the Shoup bound applies to them; the
remaining factor is a constant. Whether the parity bound of
\S\ref{sec:c25519} (at least $1.5\times$ rho with one unknown base
logarithm) applies to them is open: the cells with a validator report
($1.54$ and $1.64$) are above it, but the producer-reported $1.32$ (p36) and
$1.39$ (p42) of the scale study are below it, and no source states the
collectors' number of unknown base logarithms or reconciles the two.

\section{Cofactor-cleared bases}\label{sec:pg-cofactor}

On prime curves with cofactor $h$, a small-abscissa base wastes a fraction
$1-1/h$ of its trials on points outside the prime subgroup. A preregistered
A/B test of 3 November on the census V2 prime curves with $h=32$
(\art{research/ic_cofactor_clear_ab_20261103/}; 448 matched pairs on 112
workloads per comparison, cloud VM graded L1, counted metrics only) compared
the raw base with \texttt{prime-projected-abscissa}, which clears the
cofactor. The pass line was met: the $S$ gain of clearing is $G_0=6.180$
$[5.915,6.456]$, $G_{0.1}=6.692$ and $G_1=9.087$, and the $h=1$ negative
control gives $0.989$. The positive control, rho-plain over rho-neg with a
pre-registered lower limit above $1.2$, failed on the V2 session at
$1.227$ $[1.149,1.311]$. By the frozen decision rule the verdict is
\texttt{invalid\_controls}, and the producer declined to reinterpret it.
Descriptively, trials per relation fall from 5{,}433 to 162 ($33.6\times$
against $h=32$), and the cleared base is $9.80\times$ $[9.33,10.30]$ matched
rho at $w=0$ against $60.6\times$ for the raw base. The study is therefore
not cited as a win. The same lever on random binary curves gave
$1.29$, $1.58$ and $3.36\times$ at $h=2,2,6$ (atlas
\texttt{LEAD-ICB-cofactor-projected-ladder-base}), and the Curve25519
miniatures of \S\ref{sec:c25519} use the projected base as their baseline.

\section{Engineering record}\label{sec:pg-engineering}

\paragraph{Base size.} The 20-bit baseline at $N=45$ (Chapter~\ref{ch:distance})
was tested against $N=64$ on three new curves, 576 matched pairs, on an AMD
L2 host (\art{research/p20_large_n_20261029/}; 4{,}704 records verified,
4{,}032 measured replays identical). The ratio $S_{64}/S_{45}$ is $1.001$
$[0.979,1.022]$ at $w=0$, $0.967$ at $w=0.1$ and $0.834$ at $w=1$:
unresolved inside the band, so $N=45$ stays. It is not a relabelling,
since the total does not fall at $w=0$ and rise at $w=1$. The run also
replicated the 32-to-45 step on new curves: $N=45$ costs $0.842$, $0.822$
and $0.751$ of $N=32$, within a few percent of the original $0.818$,
$0.798$, $0.726$.

\paragraph{Where the cost goes.} A phase profile on the deployed lab image
(\art{research/ic_profile_20261029/}; 20-bit curve, $N=45$, admitted at L2,
144 of 144 executions) puts 71.8\% of priced GAE in relation collection,
27.8\% in oracle setup and nothing in the factor base, with $S_{\rm IC}=9.93$
against $S_\rho=1.08$. Amortising setup over $P$ targets gives at most a 28\%
gain ($9.93\to7.17$ as $P\to\infty$) and leaves the pipeline at $6.7\times$
matched rho; the relation phase alone is $4.7\times$ rho. To reach parity
the relation phase would have to fall by about $7\times$, which no setup
lever touches. \texttt{target\_guard\_probes} is uncharged in that record,
so $9.93$ is a lower bound.

\paragraph{Table build and membership filter.} Three builders of the same
negation-folded pair table (scalar hash insert, one batched inversion per
row, and the batched arithmetic inserted into cache-sized packed regions)
count the same additions and recover the same logarithms on a 32-bit
curve, 456 of 456 executions identical, so $S$ is unchanged by
construction; the wall comparison is L1 and descriptive
(\art{research/ic_oracle_setup_blocked_20261020/}; leaderboard component
section). A one-bit membership filter on the packed table, frozen in the
backlog item \texttt{XP-20261029-packed-filter}, is the only validated
component of the generic track (\art{benchmarks/20261029-ic-packed-filter/VALIDATION.json};
validated with caveats, single session, not replicated; EKS L2 on an AMD
EPYC 9R45). At 32 bits it lowers whole-solve wall per charged GAE to $0.9601$
$[0.9569,0.9633]$, $-3.99\%$, at identical charged GAE. In $S$ it decides
nothing: $1.0$, $1.0505$ and $1.3586$ at $w=0,0.1,1$, so the filtered table
is costlier at $w>0$. At 24 bits it reads a regression ($1.0032$) and at 28
bits it is undecidable. The distance remaining to rho-neg is $31.47$--$31.63\times$.
The validator re-priced the bit operations with its own code rather than
re-running \texttt{ecbench work}. The filter is a component speed-up outside
tables A--C of the leaderboard and moves no ranked row.

\paragraph{A fixed-window generator.} In crypto-autoresearcher, a
fixed-window generator of the $aP+bQ$ targets cut complete-solve group
operations from 200{,}856 to 50{,}490 summed over 54 toy solves per arm, with
a paired full-solve reference/candidate ratio of $1.161$ $[1.127,1.198]$ and
an A/A of $1.044$ (\art{crypto-autoresearcher:ledger/evidence/EV-ICPERF-92c130.yaml};
nine public curves at 12--16 bits, one machine and one session). Five of the
54 pairs are slower, and the record states that the result does not
replicate across machines, runtimes or sessions; the historical generator
stays the default.

\paragraph{Curve generation.} Every \texttt{prime\_search} workload had
been built by counting points with a Legendre symbol at every abscissa, about
$p$ symbols per candidate: 11\,s per 24-bit and 218\,s per 26-bit curve, paid
again by every plan, run, verify and CI replay. Stepping one point across the
Hasse interval, about $4\sqrt p$ affine additions, answers the only question
the search asks, whether $\#E$ is prime; it draws no randomness, so it
returns the same curves. Planning time fell by $5\times$ at 16 bits,
$239\times$ at 24 and $1{,}592\times$ at 26
(\art{research/prime_order_search_20261017/}). This is tooling; it changes no
counted cost.

\paragraph{Line layout.} A cache-line layout tournament of the pair table
(\art{research/ic_line_layout_tournament_20261031/}) has not produced a
measurement. The two preregistered EKS sessions of version~1 passed their
contract and build jobs and terminated in the experiment job with
remote-result exit code 23 (\texttt{TERMINAL-OBSERVATION-20261101.md}); the
version~5 qualification job of 6 November stopped at its qualification
under stopping rule~1, with no scientific session and no seed consumed. A
version~6 redesign was approved on 6 November (\texttt{DEC-20261106-662e33})
conditional on a pre-freeze check, and nothing has been dispatched under
it. The outcome is inconclusive, and the failure is tracked as
\art{problems/ic-line-layout-eks-failure-cause.md}.

\paragraph{Planner weight rules.} Chapter~\ref{sec:prime} reports the
planner's validation on six curves. The re-test of 26 October
(\art{research/ic_planner_exhaustive_20261026/}) held out three weight rules
and withdrew all three in \texttt{ic-planner/4}: the prime rule W1
($1.41\cdot N_0$ at $w=1$) held its operative half on 4 of 4 curves but
failed its $w=0$ control half, because at 15.7 and 17.7 bits $1.41N_0$ is
cheaper even at $w=0$; the binary rule W2 failed at $w=1$; the Koblitz rule
W3 was dearer at every weight on every held-out curve at 36.6--47.2 bits
($1.06$--$1.17\times$ at $w=1$). \texttt{ic plan --native-weight} now keeps
the counted optimum at every weight.

\section{Structured bases in measurement}\label{sec:pg-structured}

The conservation results of Chapter~\ref{ch:cover}
(Proposition~\ref{prop:conservation}) and Chapter~\ref{sec:transfer}
(Lemma~\ref{lem:conservation}) say that a factor base chosen by a
coordinate rule cannot change the number of relations in expectation.
Three prime-field measurements of October tested the finite-size version.

\paragraph{A four-family comparison.} Multiplicative-coset, low-$x$,
low-digit-weight and random bases were compared on 48 admitted prime-order
curves over 24 fields at 16, 20 and 24 bits, holding the target curve fixed
and matching the number of $\pm$ point classes
(\art{research/factor_base_evaluation_20261015/}; 26{,}497 point-relation
identities and 402 matrix ranks checked in Sage). In independent rows per
CPU second at 24 bits, low-$x$ gave 1{,}352, coset 1{,}167, random 1{,}144 and
low digit weight 859. Coset over random is $1.020\times$ with a descriptive
95\% interval of $0.959$--$1.083$, so no coset advantage is established. A
quadratic-algebra norm solver produced 90 rows in 229.902 CPU seconds on
fresh 24-bit inputs, against 115 rows in 0.268 seconds for the exact join;
the norm solver is a Sage prototype and the join native C.

\paragraph{Pair-sum coincidences.} A structured base might win if its
pair-sum coincidence excess $\kappa$, the number of four-term relations
$P+Q=R+S$ among base points over that of matched random symmetric bases,
grew with the base size $N$. A registered break-even ladder of 29 October
(\art{research/prime_coincidence_harvest_20261029/}) measured window,
rank-2 and rank-3 label-progression bases on prime-order curves of 16 to 32
bits at two base-size exponents. The slopes $d\ln\kappa/d\ln N$ have 95\%
intervals inside $[-0.006,+0.008]$ for all six combinations, $18$--$41\times$
below the break-even exponents $1/3$, $1/5$ and $1/7$, and every rung's
$\kappa$ lies inside $[0.92,1.08]$. The $\ZZ/N$ positive control recovers
slopes $0.699$ and $0.649$ (exact $0.70$ and $0.65$). A growing harvest would
need $\kappa>1.27\,cN^{1/(2h-1)}$, that is $1{,}874c$, $101c$ and $29c$ at 32
bits for $h=2,3,4$, against a measured $\kappa$ of at most $1.07$. The
ladder is registered as \texttt{exhausted\_at\_scale} at 32 bits; it is
not a proof.

\paragraph{Hashed and windowed bases.} Two of the skeptic-reviewed leads of
\S\ref{sec:pa-sat} add the second moment: hashed bases (affine, dense,
sparse of weight at least 5) match a size-matched random null to
$1.000\pm0.0003$ in paired ratio at $p\approx2^{14}$--$2^{22}$, and a shared
negation-closed $x$-window has a yield deficit of $0.713$, $0.833$ and
$0.916$ at $m=4$, an instance of the symmetric-base mechanism already
recorded in the knowledge base. Base choice moves constants, in the
direction the conservation results predict, and the structured P-256
bases of Chapter~\ref{sec:prime} are the same finding at a larger
modulus.

\section{The obligation chain}\label{sec:pg-obligations}

The weak-assumption programme of 2 November
(\art{research/ic_weak_assumptions_20261102/}; reanalysis of committed
evidence, no new run) restates an index-calculus security argument as a
serial chain of six obligations, each of which a complexity claim must
discharge:
\begin{itemize}
\item solution geometry: dimension and radicality of the full system,
  rational lifts and exceptional charts;
\item solver growth: matrix dimensions, nonzeros, arithmetic and memory,
  with timeouts and out-of-memory exits retained;
\item yield and variance: rational decompositions per natural target,
  including zero mass, multiplicity and solver failures;
\item relation rank: rank growth, dead columns and full-rank probability,
  not raw row counts;
\item a non-generic interface: a coordinate predicate whose useful
  candidate set no generic table reproduces at the same cost;
\item end-to-end composition: one parameter family charging setup, failed
  queries, collection, verification, linear algebra and memory, with the
  fixed-arity gates $a<m/2$ and $e<m/2-1$ of \S\ref{sec:pg-gate}.
\end{itemize}
The machine-readable chain (\texttt{obligations.json}, checked by
\texttt{audit.py}) covers naive Semaev, generic pair tables, the Kummer
tower, the proposed arity-three hybrid solver, Wagner filters, the
hashed-base tensor route of \S\ref{sec:pt-tensor} and xedni. Three of its
reanalyses bear on the generic track. The prime $m=2$ yield sweep tracks the
size-driven $2B^2/r$ model at toy scale, with observed/model hit-rate
ratios $0.764$--$1.108$ in nine cells; exact double counting fixes a mean
decomposition count, not the probability of at least one, and the
adaptively stopped collector cannot estimate the target-wise variance. The
same records hold 842 rows at rank 842, a correctness control and not a
rank theorem, since each episode stops early and meet-in-the-middle and
subtraction expose the same relation set. Its observation about degree
alone is in \S\ref{sec:pa-pkm}.

\paragraph{What the generic track settles.} On every prime curve measured
the generic track is behind rho with negation, and the best arm, the
equivariant residual collector, is behind it by a constant at rho's
exponent. The $k$-sum tables improve the incumbent by $1.4$--$4.6\times$
at 20--32 bits and stay $5.4$--$11.0\times$ rho, above baby-step giant-step.
The pair-table constants (base size, filter, builders, profile) are
measured and small. The obstruction is a theorem, not a measurement: any
decomposition oracle built from enumeration or table lookup is generic, so
Shoup's $T\gtrsim\sqrt N$ and, with preprocessing, $ST^2\gtrsim N$ bound it
(atlas \texttt{LEAD-ICP-generic-table-decomposition}, closed\_obstructed).
Every improvement here is a success of the generic track under the
iterative criterion and a distance remaining of at least $1.3\times$, and
none of it is evidence for the algebraic track.

\chapter{Prime fields: algebraic presentations}\label{ch:primealgebraic}

The algebraic track holds relations produced by a solver that receives the
factor-base constraint inside its system. Its prime-field baseline was
empty when the track was created on 29 October: the owner accepted the
small-abscissa base with membership polynomial $F_V=\prod_{c\in V}(x-c)$ as
an algebraically defined base, which makes the naive Semaev system with
msolve the counted algebraic reference, but noted that $\deg F_V=\abs V$
gives a Gr\"obner solver no low-degree structure
(\art{research/prime_ic_program_20261029/DECISIONS.md}, D1). The reference arm was approved for EKS Spot (decision D-EKS), and the
decisions file lists its implementation and execution as not yet done. The end-to-end prime-field costs nearest to that track in the treatise
are the Kummer-tower rows of Chapter~\ref{sec:prime} ($28$--$159\times$ rho),
which Chapter~\ref{ch:distance} places on the algebraic track although their
oracles are root finding and tables, so that the tower base is the only
algebraic ingredient; the P-224-shaped rows and the derived F4 cost of
\S\ref{sec:pa-pkm}; and the F4 rows of the Curve25519 miniatures
(\S\ref{sec:c25519}).

This chapter collects seven lines of work, none of which produced a priced
pipeline: the Kummer-tower follow-ups (\S\ref{sec:pa-pkm}), the partitioned
summation polynomials of Amadori, Pintore and Sala (\S\ref{sec:pa-aps}), the
digit presentation (\S\ref{sec:pa-digit}), split-graph and secant
presentations (\S\ref{sec:pa-secant}), the divisor-norm direct index calculus
(\S\ref{sec:pa-divisor}), the SAT and CDCL(T) lane (\S\ref{sec:pa-sat}) and
the small-root reach theorems (\S\ref{sec:pa-smallroot}). Each is a degree
study, a correctness implementation or a theorem. The first fall degree and
its relatives are defined in Chapter~\ref{ch:falldegree}, and that
chapter's \S\ref{sec:fd-pkm} gives the Kummer-tower degree theorem
(Theorem~\ref{thm:pkm}). None of the studies below has a verified
\texttt{pcc} claim; where a claim id exists it is named with its standing.
The fixed-arity gate of \S\ref{sec:pg-gate}, $a<m/2$ and $e<m/2-1$, applies
to every route here.

\section{Kummer-tower follow-ups}\label{sec:pa-pkm}

\paragraph{The symmetric tower.} Petit, Kosters and
Messeng's~\cite{pkm2016} base $V=g\mu_{2^t}$ makes membership of distinct
summands $x_i=gw_i$ equivalent to $Q(X)=\prod(X-w_i)$ dividing
$X^{2^t}-1$, that is $X^{2^t}\equiv1\pmod Q$. The symmetric tower of 26
October (\art{research/pkm_symmetric_tower_20261026/}; stage diagnostic,
no $S$) takes the coefficients $e_k$ of $Q$ and of the remainders
$R_j\equiv X^{2^j}\bmod Q$ as unknowns, with equations $R_j=R_{j-1}^2\bmod Q$,
$R_{t-1}^2\equiv1$ and $\Sem{m+1}$ rewritten in the $e_k$. It sees unordered
sets of distinct points, the same relation with $m!$ fewer solutions.
Table~\ref{tab:pa-sym} gives the cells. At $m=3$ the symmetric tower is one
degree lower and its widest matrix about $10\times$ narrower, about
$50\times$ faster per system at $N=12$ ($14.8$\,s against $0.3$\,s), but its
degree still climbs: $4,5,5,6$ at $N=6,9,12,15$. The follow-up of the same
day (\art{research/pkm_symmetric_lead1_20261026/}) added $N=18$ at $m=3$,
where $D=6$ again, so the sequence over $N=6$--$18$ reads $4,5,5,6,6$: a
staircase that steps up every second cell. The $N=18$ cell cannot tell
``$D$ bounded at 6'' from ``$D$ rises by one every 6 in $N$''; the
discriminating cell is $N=21$, which needs more than the 13\,GB that host
had. From $N=15$ to 18 the widest matrix grows $3.1\times$ and the time
$39\times$. At $m=4$, $\Sem5$ was written in $e_1,\dots,e_4$ by
interpolation and passed its planted check. At $N=12$ the symmetric system
runs in 8--9\,ms against 66--70\,s for the plain chain, with
$5.9\times$ narrower matrices; its $D=8$ at $N=12$ and 16 is the degree of
$\Sem5(e)$, its input, so $D$ cannot rank the arms at $m=4$, and $N=20$
aborted under a 6\,GB cap.

\begin{table}[t]
\centering\small
\begin{tabular}{@{}rrrrrr@{}}
\toprule
$m$ & $N$ & plain $D$ & symmetric $D$ & widest matrix, plain / sym & time per system, plain / sym \\
\midrule
2 & 12 & 5 & 4 & 2{,}451 / 627 & 0.36\,s / 0.09\,s \\
2 & 16 & 5 & 5 & 9{,}119 / 6{,}281 & 34\,s / 24\,s \\
2 & 18 & 5 & 5 & 14{,}922 / 12{,}214 & 280\,s / 176\,s \\
3 & 9 & 6 & 5 & 3{,}285 / 248 & 0.6\,s / 0.02\,s \\
3 & 12 & 6 & 5 & 12{,}314 / 1{,}224 & 14.8\,s / 0.3\,s \\
3 & 15 & $\ge6$ & 6 & $\ge97{,}456$ / 10{,}532 & $>900$\,s / 43\,s \\
3 & 18 & --- & 6 & --- / 33{,}153 & --- / 1{,}941\,s \\
4 & 12 & 7 & 8 & 21{,}619 / 3{,}663 & 66--70\,s / 8--9\,ms \\
4 & 16 & --- & 8 & --- / 13{,}926 & --- / 11.8\,s \\
\bottomrule
\end{tabular}
\caption[]{Plain and symmetric Kummer towers, F4 (\texttt{f4\_fp}, grevlex)
on $p=786433$, random targets; maximum over two targets. Stage diagnostics
with no priced pipeline. The recorded tower-engine $D$ of the plain tower at
$m=3$, $N=15$ is 7 (\S\ref{sec:pkm}). Sources:
\art{research/pkm_symmetric_tower_20261026/},
\art{research/pkm_symmetric_lead1_20261026/}.}
\label{tab:pa-sym}
\end{table}

\paragraph{What the tower instrument does not check.} The 29 October audit
(\art{research/prime_one_shot_20261029/}) read the source of
\texttt{examples/pkm\_symmetric\_tower.rs} and found that it draws $a,b$
without checking the discriminant or the group order, includes every coset
abscissa in the base without checking that it lifts to $E(\FF_p)$, and
takes $x_R$ as a uniform field element rather than the abscissa of a
subgroup point. Its planted checks evaluate assignments built from $\Sem3$
roots, not signed group decompositions. The symmetric rows are therefore
algebraic instrument diagnostics, not a validated natural prime-order
benchmark, and the symmetric arm excludes repeated abscissae that the plain
arm permits, so the comparison is not clean evidence for coefficient
sharing alone. The same audit recovered a lower bound the earlier
summaries had dropped: at $m=4$, $N=20$ both raw msolve logs complete three
productive degree-8 rounds (adding 1{,}241, 2{,}507 and 3{,}578 basis
elements) before exit 139, so $D\ge8$ there for that solver and
presentation, whatever the final verdict.

The audit sorted the tower's failures and nulls into kinds that need
different treatment, and its sorting applies beyond the tower:
\begin{itemize}
\item \emph{Busy hosts.} The original PKM solver study ran at load 22--25 on
  four cores; a single-threaded $m=3$ job received about 19\% of one core,
  so the first degree-7 steps at $N=18$ and 21 took about 2{,}024 and
  2{,}041\,s wall against 384 and 396\,s CPU. The wall-censored inputs
  justify a controlled retry, not a reading of the exponent. The same study
  planned a planted-solution control but retained only sixteen refutations
  with \texttt{--planted 0}, so successful root extraction was not validated
  there.
\item \emph{Resource censoring.} At $m=3$, $N=21$ the native engine stopped
  before productive degree-8 reduction because a dense block would need
  79{,}330{,}205{,}871 entries, about 317\,GB at four bytes, which would not
  fit the planned 192\,GB host either. That is a matrix-memory obstruction,
  not a contention artefact, and it is why \S\ref{sec:pkm} reads $D\ge8$ as
  inconclusive.
\item \emph{Missing raw runs.} The last-fall study's raw runs, table and
  draft results were never pushed and live only in the owner's laptop
  worktree (\art{research/prime_fall_degree_20261027/RESTART.md}), so its
  quoted medians are session-reported and not independently recoverable. A
  lost file is not a negative experiment; decision \texttt{D-MAC} asks for
  recovery.
\item \emph{Insensitive instruments.} An ecbench Latt\`es and isogeny
  study used a degree-$s$ top product polynomial, $s=16,24,32$, that set the
  solving degree by itself; its null at $m=2$, $p\le1031$ is not a theorem
  about implicit tower presentations, and rerunning the same design would
  add little.
\end{itemize}

\paragraph{Degree alone is not a cost certificate.} The weak-assumption
reanalysis (\art{research/ic_weak_assumptions_20261102/}) extracted the
committed PKM $m=4$ random-target rows at $N=8,12,16$. Their F4 solving
degrees are $6,7,7$, while the maximum column count grows $1{,}049\to9{,}452
\to76{,}837$ and multiply-adds $4.84\cdot10^7\to8.55\cdot10^{10}\to
2.70\cdot10^{13}$: from $N=8$ to 16 the columns grow $73.25$-fold and the
work $558{,}304$-fold, with the last two degrees equal. Three cells are not
an asymptotic fit, but they are a direct counterexample to reading a stable
degree as a bounded cost. All six selected $m=4$ natural rows are also
inconsistent (their ideal is the unit ideal), so a zero-dimensional flag on
a random target does not establish the geometry of a successful
decomposition.

\paragraph{P-224's 2-adic shape.} A study of 22 October priced the tower at
P-224-shaped Solinas primes $p=2^n-2^k+1$ of 20--61 bits and at the P-224
prime itself (\art{research/pkm_p224_kummer_shape_20261022/}; frozen
before any measured run). On the solver axis, $D$ equals the recorded $D$
in all 89 valued tower cells at $m=2,3,4$, the width and step count of
every system match, and the shape-matched null gives the same $D$ in all 65
paired cells, although $k/n$ ran from $0.42$ to $0.70$ and the odd part of
$p-1$ from 6 to 128 bits. The one cell past the recorded frontier,
$m=2$, $N=24$, was killed at 6.0\,GB, and a capped rerun gives only $D\ge6$.
On price, four Solinas curves of 20--28 bits cost $18$--$78\times$ rho with
negation at $w=0$ and $19.9$--$134\times$ at $w=1$, every one of 1{,}152 runs
verified and 864 replays reproduced; at equal $t$ the Solinas curves cost
2--32\% more absolute GAE than the earlier Proth curves. The derived
end-to-end F4 cost at the one cell matching a curve (20 bits, $m=2$,
$N=20$) is at least $2.4\cdot10^9$ in $S$, about $1.9\cdot10^9\times$
matched rho at $w=0$. The verdict is inconclusive on the solver axis and
negative on price.

\paragraph{Radix gradings.} The experiment \texttt{EXP-PKM-RADIX-001}
(\art{research/pkm_radix_grading_20261104/REPORT.md}; frozen protocol, two
disclosed amendments; every number is a counted metric marked
\texttt{[pcc:unverified]} in the report) regraded the tower with radix 4 and
radix 8 and measured refutation multiply-adds. Its frozen verdicts are
\texttt{gain\_constant} at $m=2$ radix 8 ($r_k=0.173$ at $t=10$), $m=3$
radix 4 ($0.344$ at $t=5$) and $m=3$ radix 8 ($0.234$), and
\texttt{no\_material\_effect} at $m=2$ radix 4 ($0.538$). No cell is
\texttt{gain\_growing}. The shape-matched null gives the same verdicts with
ratios equal to three digits, so the gain belongs to the support of $\Sem3$
on the tower, not to the curve; $p=786433$ replicates the $t=10$ bands. The
author's registered prediction of no material effect anywhere was
falsified in three of four cells. Table~\ref{tab:pa-radix} gives the end cells.

\begin{table}[t]
\centering\small
\begin{tabular}{@{}lrrrll@{}}
\toprule
Cell & radix 2 & radix 4 & radix 8 & $D$ (2/4/8) & widest matrix (2/4/8) \\
\midrule
$m=2$, $t=5$ ($N=10$) & $5.43\cdot10^6$ & $2.20\cdot10^6$ & $1.64\cdot10^6$ & 4 / 6 / 10 & 337 / 252 / 260 \\
$m=2$, $t=10$ ($N=20$) & $4.82\cdot10^{11}$ & $2.59\cdot10^{11}$ & $8.35\cdot10^{10}$ & 6 / 8 / 13 & 33{,}446 / 28{,}225 / 17{,}101 \\
$m=3$, $t=4$ ($N=12$) & $1.06\cdot10^{10}$ & $9.59\cdot10^9$ & $5.98\cdot10^9$ & 6 / 9 / 12 & 5{,}019 / 6{,}509 / 4{,}063 \\
$m=3$, $t=5$ ($N=15$) & $3.19\cdot10^{12}$ & $1.10\cdot10^{12}$ & $7.47\cdot10^{11}$ & 7 / 9 / 12 & 31{,}397 / 28{,}055 / 20{,}284 \\
\bottomrule
\end{tabular}
\caption[]{Refutation multiply-adds of the Kummer tower in three gradings,
one value per cell (the two random draws agree to a few parts in $10^7$);
$D$ is measured in each grading's own degree. Counted metrics,
\texttt{[pcc:unverified]}. Source:
\art{research/pkm_radix_grading_20261104/REPORT.md}.}
\label{tab:pa-radix}
\end{table} The radices raise the degree in their own
grading (at $m=2$, $t=10$: $6$, $8$, $13$) while lowering the work, which is
the point of the previous paragraph seen from the other side.

\paragraph{Isogeny kernels against $\mu_k$.} A component study of 2
November tested a user's toy observation that an isogeny-kernel base,
$V=x(\langle G\rangle)\subset E'(\FF_p)$ presented as a chain of V\'elu
2-isogeny $x$-maps, behaves differently from Petit's $\mu_k$ chain of
squarings at $m=2$ (\art{research/isogeny_vs_mu_factor_base_20261102/};
toy primes $p\le40961$, frozen protocols, 6{,}144 and 11{,}532 systems, no
gate failures, local host without an isolation tier, counted units only).
At equal $\abs V$ the isogeny base matches a random base in yield in 15 of 15
cells (median difference $-0.020$ to $+0.029$). Its extra Gr\"obner cost
comes from the presentation, not the set: writing $\mu_k$ in coordinates
whose steps look like V\'elu steps costs the same as the isogeny chain to
within 3\%. A search over coordinates cut the lex-order gap to $\mu_k$ from
$6.4$--$11.7\times$ to $1.9$--$2.6\times$, but found no grevlex improvement
(0 of 16 cells), and none of 11{,}811 valid V\'elu representations reaches
the density of the $\mu_k$ chain. There is no rho row; the distance is
stated against $\mu_k$.

\paragraph{The $t=14$ cell.} The decisive cell of the prime fall-degree
study (Kummer, $m=2$, $p=2013265921$, $t=14$) is half done
(\art{research/prime_fall_degree_t14_20261101/}). At cap 5 neither random
draw refutes and neither is decomposable, so $d_F\ge6$. The cap-6 cells
need a 112\,GiB, 86{,}400\,s native job; no route in the lab carried such a
job on 1 November, and the cells are owed. The owner approved a command
route for non-ecbench jobs (\texttt{D-ROUTE}) that an administrator has to
implement.

\paragraph{Open.} Lead~1 of the tower, whether $D$ stays bounded as $N$
grows, is open in every presentation tried: plain, torus, Dickson,
symmetric, radix-graded and P-224-shaped. In each family $D$ either grows
with $N$ or is pinned by the input degree, and the shape-matched null
matches. The atlas lead \texttt{LEAD-ICP-pkm-tower-algebraic-oracle} is
\texttt{open\_blocked} on memory: the cells that would discriminate
($m=3$, $N=21$ in the symmetric tower; $t=14$ at cap 6) need hosts the
programme did not have. Chapter~\ref{ch:falldegree} gives the proved part.

\section{Partitioned summation polynomials}\label{sec:pa-aps}

Amadori, Pintore and Sala~\cite{aps2018} partition the factor base and
attach labels to the summands so that one relation comes from a smaller
system. The programme implemented the paper's single-relation route and
five extensions between 26 October and 5 November
(\art{research/aps_extensions_20261026/},
\art{research/aps_followup_20261028/},
\art{research/aps_independent_replay_20261029/},
\art{research/aps_measurement_20261029/},
\art{research/aps_export_20261104/},
\art{research/aps_leads_20261105/}). The native module builds labelled
points $R=uG+vQ$ without giving the solver the planted scalar, imposes the
partition membership constraints, solves with native F4 and substitution,
decodes and sign-checks the summands by curve addition, and recovers
$d=-A/B$ from accumulated relations $AG+BQ=O$. The extensions are fixed
coefficient labels, unordered divisor and two-pair encodings, automorphism
orbits, covering partitions, and expanded, split and hybrid solving; the
follow-up added an early stop on the first informative relation. In
detail:
\begin{itemize}
\item \emph{Informative labels.} For odd arity, points $uG+Q$ have an odd
  signed coefficient sum, which cannot vanish modulo $n>m$; for even arity,
  one point from each partition takes weights $1,2,4,\dots$ with
  $n>\sum$weights. Regrouping or orbit expansion does not inherit the
  guarantee, so the recovery check stays mandatory.
\item \emph{Unordered selections.} A monic cubic
  $H(T)=T^3-e_1T^2+e_2T-e_3$ dividing $f(T)=\prod_{v\in V}(T-v)$ selects three
  distinct base coordinates exactly when the three remainder coefficients
  vanish, with the symmetric $\Sem3$ rewrite $(e_2-a)^2-4e_1(e_3+b)=0$; the
  study notes that the remainder equations can be much harder to solve and
  claims no degree reduction.
\item \emph{Automorphism orbits.} On a $j=0$ fixture the map
  $\varphi(x,y)=(\beta x,y)$ propagates labels through orbits; the intrinsic
  relation $R+\varphi(R)+\varphi^2(R)=O$ has $A=B=0$ and is rejected.
\item \emph{Covering partitions.} A bit-based three-colour separating
  family covers every distinct triple with at most
  $2\lceil\log_2s\rceil(\lceil\log_2s\rceil-1)$ colourings, an $O(\log^2s)$
  construction; coverage does not imply a lower total cost.
\item \emph{Solver presentations.} Expanded $\Sem4$, the paper's joined-$\Sem3$
  form and first-coordinate hybrid guessing share one checked relation path.
\end{itemize}
The field is bounded by $65{,}521$ by construction, and the study states
that success of its substitution backend on the fixtures is no evidence
that enumeration disappears at large sizes.

What exists is correctness on tiny public fixtures. The independent replay
reads the correctness JSON for three fixed fixtures and twelve arms and
re-checks group order, every labelled point, partition membership, signed
sums, $A$, $B$ and $dG=Q$ in separate affine arithmetic. The follow-up's
release suite reports 5{,}002 tests passed, 0 failed and 122 ignored, with
seven core-only Lean lemmas. The measurement instrument \texttt{ic.aps\_v1}
charges setup and five online phases on one monotonic cold timeline, but
fields are capped at 65{,}521, abscissae at 32 and attempts at 128; its
\texttt{total\_gae} prices only the independent affine checks, so it is a
lower bound and \texttt{native\_work\_accounting\_complete} stays false. No
benchmark has launched, and the paper's 11--22-bit ladder was not
reproduced. The obstruction record
\art{problems/aps-single-relation-complete-cost.md} lists what a complete
cost needs. The four leads of 5 November (a nonhorizontal line chart
$U=t^2$, $V=c/t$ with $H_{U,V}(X)=X^3-U(X+V)^2+aX+b$; F4 trace reuse across
attempts; incremental factor-base frontiers; unequal partition profiles)
are hypotheses with no implementation, filed as \texttt{PCC-20261105-aca921}
and \texttt{PCC-20261105-ac1e93}, both of unverified standing. The APS
route is open and unmeasured.

\section{The digit presentation}\label{sec:pa-digit}

Writing each unknown abscissa in $s$ bounded binary digits,
$\ell_k=\sum_{i<s}2^ia_{k,i}$ with $a_{k,i}^2=a_{k,i}$, gives a prime-field
analogue of Weil descent: $n=ms$ Boolean unknowns, and the system
$\tilde S=\Sem{m+1}(\ell_1,\dots,\ell_m,x_R)$ in the function algebra
$\FF_p[a]/(a^2-a)$. The closed form of its first fall degree,
$d_{\rm ff}=m2^{m-1}+\lfloor(s-2^{m-1})/2\rfloor+1$, is in
Chapter~\ref{ch:summer}, and its $m=2$ case $4+\lfloor s/2\rfloor$ in
Chapter~\ref{ch:falldegree}. A native $\FF_p$ ladder with an independent
prime, curve and target for every draw
(\art{research/prime_digit_presentation_20261029/}; registered before the
counted runs, one amendment) measured the \emph{last} fall degree:
\begin{itemize}
\item $d_{\rm last}=4+\lfloor s/2\rfloor$ on 106 of 106 real draws at $m=2$
  (60 unsatisfiable at $s=2$--7, 46 planted at $s=2$--8), equal to the first
  fall degree in every draw; at $m=3$, the first fall degree is 13 on 8
  independent instances at $s=4,5$, and $d_{\rm last}=13$ on the 4 at $s=4$;
\item a support-matched null gives $s+3$ on 44 of 44 draws, and the real
  system falls strictly earlier from $s=3$; a block-factored null of random
  forms on the support of $\ell_k^e$ reproduces the real first and last fall
  degrees on 44 of 44 draws, so the early fall is block structure and
  carries nothing curve-specific;
\item the plain presentation $\FF_p[x_1,x_2]/(x_i^p-x_i)$ in the same engine
  has $d_{\rm last}=p+2$ on 32 of 32 draws: linear in $p$, where the digit
  degree is linear in $\log p$;
\item at $d_{\rm last}$ the degree-filtered Macaulay space spans 100, 98,
  96, 83, 81 and 60\% of the $2^n$ monomials at $s=2$--7, and 60\% at $s=8$.
\end{itemize}
\begin{table}[t]
\centering\small
\begin{tabular}{@{}rrrrrrr@{}}
\toprule
$s$ & draws & distinct primes & $d_{\rm last}=d_{\rm ff}$ & support-matched null & columns at $d_{\rm last}$ & columns$/2^n$ \\
\midrule
2 & 12 & 5 & 5 & 5 & 16 & 1.00 \\
3 & 12 & 10 & 5 & 6 & 63 & 0.98 \\
4 & 12 & 11 & 6 & 7 & 247 & 0.96 \\
5 & 12 & 11 & 6 & 8 & 848 & 0.83 \\
6 & 8 & 8 & 7 & 9 & 3{,}302 & 0.81 \\
7 & 4 & 4 & 7 & 10 & 9{,}908 & 0.60 \\
\bottomrule
\end{tabular}
\caption[]{Digit-presented Semaev system at $m=2$, unsatisfiable real
draws, each with its own prime, curve and target; none censored.
``Columns'' is $\sum_{j\le d_{\rm last}}\binom nj$, derived from the
measured degree. Source: \art{research/prime_digit_presentation_20261029/}.}
\label{tab:pa-digit}
\end{table}

Table~\ref{tab:pa-digit} gives the unsatisfiable ladder. A low degree does not make a small system at this scale, and the lane is
exhausted at the measured scale ($m=2$, $s\le8$; $m=3$, $s=4,5$). It would
reopen if $s=9$ broke the closed form (predicted $d_{\rm last}=8$, about
$2^{18}$ columns) or if the adaptive staircase arm, not run, fell below the
block null. The scale limit is the engine's \texttt{u32} accumulators
($p<2^{12}$), not mathematics; for $s\ge8$ no unsatisfiable target remains on
the $2^{16}$-point grid, so $s=8$ has only two planted draws ($d_{\rm last}=8$,
39{,}203 columns, about 11 minutes a draw). Validation draws preceded
registration, and the registered binary's hash is not reproduced by a
rebuild, although the rebuilt source reproduces two of panel~A's cells line
for line. No new theorem is claimed: \texttt{lean/NullTable.lean} checks
the registered table's arithmetic, and a faithful proof of the closed form
is blocked on Wilson's rank theorem for inclusion matrices and two further
dependencies absent from core Lean. A frozen but unrun protocol for a $p^3$
second-digit carry experiment is in
\art{research/p3_second_digit_carry_20261101/}.

\section{Split-graph and secant presentations}\label{sec:pa-secant}

Take a factor base of the rational intersections of $E$ with the graph of a
low-degree ordinate function $y=h(x)$ whose $h^2-f$ splits completely, and
write a two-point decomposition through the secant through the two
summands. The split-graph theorem of 31 October
(\art{research/prime_factor_base_degree_20261031/SECANT_GRAPH.md}) proves a
homogeneous top-ideal degree $d+1$ for this presentation against $2d+1$ for
the same base in the dense two-summand system, and a Lean continuation
proves the polynomial-ring reduction with the actual highest-degree forms,
including the exact $\FF_{17}$ fixture at degree $5\to3$ for every target and
all four sign branches (159 declarations in fourteen modules, zero failures,
no native trust). The measured part
(\art{research/prime_factor_base_degree_20261031/parallel_20261031/RESULTS.md})
ran stock F4 on 14{,}220 systems, of which 14{,}194 verified and 26 hit
root-extraction caps (all at $\FF_{127}$, degree 5). The handled F4 degree of
the secant presentation is 3 against 6 for the dense presentation on all 18
$\FF_{17}$ targets, and $d+1$ against $2d+2$ on $\FF_{83}$, $\FF_{101}$ and
$\FF_{127}$ at graph degree $d=3$--5, with the matched non-graph secant
control in between at $2d$.

The degree gain does not become a cost gain. Against the matched secant
control the candidate improves degree, width and pair and row counts.
Against its own dense presentation it lowers the degree but increases the
branch-summed work: on the $\FF_{17}$ toy, 976 reduced pairs and 2{,}102
reduced rows against 344 and 736 for the dense system. The stock engine
exposes no field-operation total, so cost per relation is null. Equal-size
bases have the same all-target ordered yield $K^2-K$. A Boolean-offset
successor of 1 November
(\art{research/prime_factor_base_degree_20261031/pencil_search_20261101/})
removes the complete-splitting requirement at canonical degree 7 against a
same-base dense 9, and selector recodings prove $7\to6$ (XOR) and $7\to5$
(one-hot); affine degree, yield and work are unmeasured. The atlas lead
\texttt{LEAD-ICP-split-graph-secant-presentation} is \texttt{open\_next\_step}:
a degree theorem with no cost advantage.

\section{Divisor-norm decompositions and a direct index calculus}\label{sec:pa-divisor}

An $m$-point relation can be written as a norm: a function in
$L((m+1)O)$ of the form $A(X)-B(X)y$ whose norm $A^2-B^2f$ is divisible by
$X-x_R$ and whose other roots must lie in the factor base. Membership in a
Kummer base $g\mu_d$ can then be imposed by quotient remainders or by
Graeffe coefficient towers. The ``one research bet'' study of 29 October
chose growing-arity divisor-coefficient decomposition with shared quotient
and remainder membership towers, and stated that it could not deliver an
original sub-exponential algorithm because the essential complexity lemma
is unproved (\art{research/prime_one_shot_20261029/}). Its continuation
(\art{research/prime_one_shot_20261029/continuation/}) built the pieces:
\begin{itemize}
\item a native coefficient solver for arity three that receives only the
  curve, target and base parameters, searches the two free norm
  coefficients, enforces membership through quotient remainders and extracts
  every rational selected-chart solution; on $E:y^2=x^3+2x+2$ over
  $\FF_{17}$ its complete model sets agree with a separate signed-point
  oracle in all 90 cases (42 SAT, 48 UNSAT, 400 models);
\item a mixed-target direct index calculus that recovers the logarithm of
  every one of the curve's 18 nonzero targets from $R=aP+bQ$, with a
  left-null vector of the factor-base coefficient matrix alone and no
  supplied logarithms, retaining 230 attempts with all 217 nonidentity
  solver calls charged (59{,}024 coefficient pairs);
\item written proofs, re-derived within the session but not formally
  certified, of rationality, a coverage bound $\rho_g\ge1/(192m16^m)$ for a
  constant fraction of admissible bases under stated finite conditions
  ($1/(320m16^m)$ for the conservative selected chart), and recovery
  without full base rank.
\end{itemize}
Each current solver call visits $p(p-1)$ coefficient pairs, so even the
coefficient loop is $\Theta(2^{2n})$ in $n=\log_2p$: the solver is
exponential. The supporting reduction would need $2^{4\sqrt n+O(\log n)}$
oracle calls and dense linear algebra of $2^{6\sqrt n+O(\log n)}$ bit
operations, so a proved uniform oracle bound $C(n)=\exp(o(n))$ would give a
sub-exponential algorithm in that field family. No such bound is supplied.
The obstruction record is \art{problems/prime-divisor-uniform-solver.md},
the atlas lead \texttt{LEAD-ICP-shared-coefficient-membership} is
\texttt{open\_next\_step}, and no exponent was measured. The same
continuation found the external tensor-contraction preprint of
\S\ref{sec:pt-tensor} and checked that the literal shared-remainder system
translates to $O(mb^2)$ Boolean edges, too many for the contraction bound
to give $\exp(o(b))$.

\section{SAT, SMT and CDCL(T)}\label{sec:pa-sat}

Chapter~\ref{sec:prime} reports three frozen SMT formulations of a 20-bit
decomposition that returned \textsc{unknown} at 120\,s. The lead review of
31 October (\art{research/prime_ic_lead_review_20261031/}) asked where
SAT-shaped methods belong. It ranked the hashed-base tensor route of
\S\ref{sec:pt-tensor} first, as the only inspected route with an explicit
claimed sub-exponential bound, and noted that this is not evidence for a
CDCL solver: the source keeps an $O(mb)$-edge tensor network because an
ordinary circuit-to-CNF reduction loses the size bound. For the narrower
practical-solver question it drafted a fixed-arity-three CDCL(T)/$\FF_p$
screen on a one-hot, disjoint-window, chained-$\Sem3$ predicate, with
strict gates $e<0.5$ and $a<1.5$, deliberately not frozen until its
generator, theory propagator and preflight cost exist. The shared-coefficient
discriminator (\art{research/prime_shared_coefficient_sat_20261031/}) passed
its E0 gate (18 of 18 exact model-set matches on $\FF_{17}$; 6 SAT and 12
UNSAT targets) and, on 1 November, its Boolean E1 gate: 36 of 36
target/backend model-set comparisons and 42 of 42 independently checked
UNSAT proofs. No timing or scaling was collected; the elimination campaign
of \S\ref{sec:pt-tensor} records the screen as
``representation gate passed; performance implementation-blocked''.

\paragraph{Twelve leads.} On 31 October twelve prime-field leads, L01--L12,
were evaluated and then reviewed by skeptic agents, whose digest is
authoritative on standing (\art{research/ic_prime_eth_leads_20261031/SYNTHESIS.md};
\art{research/ic_prime_eth_leads_20261031/leads/SKEPTIC_DIGEST.md}).
Table~\ref{tab:pa-eth} gives the standing of each lead after review; the
full account of what the leads measured, the refutations and the
synthesis's feasibility limits is in \S\ref{sec:tn-leads}. The synthesis's
verdict is that nothing moved the exponent question, and that the SAT and
propagation lane for prime fields ``should be closed for exponent
purposes''; it leaves open the Kummer-tower cells at $m\ge4$
(\S\ref{sec:pa-pkm}) and the contraction exponent of the tensor route
(\S\ref{sec:pt-tensor}).

\begin{table}[t]
\centering\small
\begin{tabularx}{\linewidth}{@{}lXX@{}}
\toprule
Lead & Headline after skeptic review & Standing \\
\midrule
L01 char-2 gadget & slope about $2.58$ bits of conflicts per window bit, against a bar of $0.5$; CPU $3\times$ worse at $\ell=7$ & constant only \\
L03 sparse-parity bases & paired ratio to a size-matched null $1.000\pm0.0003$, $b\le22$ & constant only \\
L04 backdoor census & offset $2\log_2(\mathrm{cap})-1$ at $m=3$: enumeration moved into the propagator & contested; cap-bounded statement only \\
L05 policy pilot & reference slope $2.17$ $[1.93,2.38]$; renaming is null by construction & refuted as stated (the claim ``closed''); policy question untested \\
L06 contraction width & MMD+ lower bound 15--51 against enumeration $2$--$10.7$ at $b\le16$ & refuted as closed; narrow claim stands \\
L07 base second moment & hashed bases $1.000\pm0.005$ of null; shared negation-closed window $0.713/0.833/0.916$ at $m=4$ & contested; constant only \\
L08 Koblitz transfer & true margin about $1.4\times$, not $3.42\times$; binary only & refuted as stated \\
L09 BTN census & margin $-16.1$ to $-16.3$ bits, no encoding admitted by D.1 ($\lambda\ge2^{85}$ needed) & tool stands, yardstick downgraded \\
L10 Lean filings & seven gaps between the lead text and Lean, all confirmed & refuted as fileable \\
L11 multiplier encodings & pc/school $0.78$ $[0.67,0.92]$ at $b=9$ only & stands, weak \\
L12 clause retention & inc/fresh $0.025$ $[0.014,0.040]$, setup uncharged & refuted; retained setup \\
\bottomrule
\end{tabularx}
\caption[]{The 31 October prime-field leads after skeptic review (L02 was
never run). Producer-reported, from the synthesis; no lead has a validation
record. Source: \art{research/ic_prime_eth_leads_20261031/SYNTHESIS.md}.}
\label{tab:pa-eth}
\end{table}

The review corrected the lead document in eight places; the corrected
sentences are in \S\ref{sec:tn-corrections}, and none changes an exponent
argument.

\section{Small-root reach}\label{sec:pa-smallroot}

A windowed decomposition asks for a root of a block $f=\Sem3(x_1,x_2,x_3)
\bmod p$ with $\abs{x_i}<p^{\delta_i}$, the reach being
$\varepsilon=\sum\delta_i$; at $(m,w)=(4,2)$ index calculus needs
$\varepsilon=8/9$. Chapter~\ref{sec:prime} records the $\Sem5$ Coppersmith
reformulation ($\delta^*=1/80$) and the $\Sem3$ falsifier at $\beta=1/4$.
Four studies of 26--29 October turned this into theorems.

\paragraph{Correction.} The atlas lead quoted a margin of $-0.048$ at
$(4,2)$. It came from \texttt{tools/smallroot\_reach.py}, which compared
$\prod X_i^{s_i}$ with $p^{s_p}$, where $s_p$ is the exponent of $p$ in
$\det L$. The Howgrave-Graham criterion compares with $p^{t\dim L-s_p}$,
and for nested shift sets $\varepsilon_{\rm HG}=\varepsilon_{\rm tool}/n$
in the limit (\texttt{CORR-20261024-b2a799};
\art{crypto-autoresearcher:analysis/smallroot-reach/CORRECTION.md}). The
Jochemsz--May extended reach of $\Sem3$ falls from $0.8408$ to $0.2803$, and
the best margin at $m=4$ from $-0.0481$ to $-0.6086$. The note is a
computation by one agent; its pre-registered audit
\texttt{EXP-ECDLP-64fa3d} has not run.

\paragraph{Any lattice, equal bounds.} The optimal-reach study
(\art{research/smallroot_optimal_reach_20261026/}) bounded every
Coppersmith lattice built from $(p,f)^t$, the provably optimal shifts of
Ryan included, without building one: with equal per-variable bounds the
reach is at most $(n!\,\mathrm{vol}\,N)^{-1/n}$ for Newton polytope $N$,
which is $0.3293$ for $\Sem3$ in three unknowns ($\mathrm{vol}\,N=14/3$) and
$0.3536$ in two. The bound is attained in 11 of 11 exact cases. Against the
demand of $0.8889$ the margin is $-0.56$, proved.

\paragraph{Any bounds, any monomial set.} The follow-up
(\art{research/smallroot_followups_20261026/}) proved that any such lattice
on any monomial set $M$, convex or not, satisfying the Howgrave-Graham
condition needs $h_N(\delta)=\max_{v\in N}\langle\delta,v\rangle<1$. For
$\Sem3$ in three unknowns this gives $\varepsilon<3/4$ and in two
$\varepsilon<1/2$, against $8/9$. An exact check found no violation in 9{,}600
cases. A joint lattice for the three-equation $(4,2)$ system needs
$\nu_s<s$ for some $s$; the computed $\nu_s/s\ge2.08$ in the demand region
for $s\le3$ is a computation, not a proof. A third study
(\art{research/prime_eform_reach_20261029/}; derived proof, exact counts)
proved Theorem~K: for a convex $M$ and the extent $W_i$ of $f$ in
coordinate $i$, every lattice gain satisfies $\sum_i\delta_iW_i\le1$, so
$\varepsilon\le1/\min_iW_i$, which is $1/2$ for $\Sem3$, the value the
numerical search had only approached. In the symmetric $e$-variables with
full extents the same argument gives $\delta_B\le2/(m(m+1)2^{m-1})$, that is
$1/6$, $1/24$ and $1/80$ at $m=2,3,4$, so the Jochemsz--May asymptotes are
optimal there. The closure is conditional: at the exceptional
specialisation $z=0$ the extents are not full, and a review counterexample
reaches $7/41=0.1707>1/6$ for $\Sem3$. Finally, any $x$-variable lattice on
a convex set reaches at most $\delta=1/21.8$ for $\Sem4$ and $1/70.8$ for
$\Sem5$, against $1/4$ and $1/5$ needed, $5.5$ and $14.2\times$ short
(\art{research/prime_s5_small_root_leads_20261028/}, \S7c); this rests on
the written proof above, which has no Lean formalisation.

The small-root family is closed for one block, for every bound shape and
every monomial set by proof, and for convex sets at the $e$-form
asymptotes. It is open for multi-target ideals, the exceptional
specialisations and joint lattices beyond $s=3$.

\chapter{Prime fields: the tensor route and named targets}\label{ch:primetargets}

This chapter has two parts. The first is a review of an external claim:
the preprint of Akshima and coauthors, \emph{A Subexponential Time Algorithm
for Elliptic Curve Discrete Logarithm}, which claims sub-exponential
prime-field ECDLP conditional on a Boolean tensor-contraction theorem in a
second external manuscript. The programme read both, ported the prototype
and built instruments. \emph{Nothing was measured}: there is no ecbench
count, no wall time and no cost against rho for this route, and the
programme did not verify the theorem it rests on. The second part collects
the named targets not already in Chapter~\ref{sec:prime}: EcGFp5 and
EcMasFp5 over $\FF_{p^5}$, a short remainder on $\FF_{p^4}$, ECCp-131, and
the audits of standard prime curves. Xedni calculus and the JKSST audit
are in Chapter~\ref{ch:obstructions}.

\section{Boolean tensor contraction: an external claim}\label{sec:pt-tensor}

\paragraph{The claim.} Version~0.2 of the Akshima et al.\ preprint (pinned
at \texttt{bochum-crypto/subexponential-ecdl} commit \texttt{1ee44e56},
PDF SHA-256 \texttt{82e27362\ldots}) claims a classical expected-time
algorithm
\[
  2^{O(b/(\log b)^{1/4})},\qquad b=\lceil\log_2p\rceil,
\]
for prime-field ECDLP with known odd prime group order, and a fixed-cofactor
extension (\art{research/prime_one_shot_20261029/continuation/tensor-review/README.md}).
Its Theorem~4 is conditional on its Theorem~2, which invokes Appendix~D.1
of \emph{The Exponential Time Hypothesis is False} (pinned at
\texttt{openai/math} commit \texttt{7ef5c0cc}, PDF SHA-256
\texttt{5b1612dc\ldots}; the entry-cost audit of 4 November found the
publisher path's later commit \texttt{42907c9e} serves the same PDF). That
appendix states a contraction bound of
$(s+L_{\rm in})^{O(1)}2^{O(s/\sqrt{\log s})}$ deterministic bit work, in
polynomial space, on an $s$-edge bipartite Boolean tensor network whose
vertex predicates are charged explicitly. The source triage records that
repository and author labels establish provenance only, not endorsement,
peer review or correctness.

\paragraph{The mechanism.} The preprint encodes integer multiplication with
$O(b)$ Boolean edges and maximum vertex degree $O(b^{3/4})$, through two
stages of exact Fourier transforms over a finite ring; a straightforward bit
circuit's extra $\log b$ factor would destroy the saving. Independent affine
hashes of injectively compressed points define $m$ separate public factor
bases, each membership costing $O(b)$ edges, and a complete $m$-point sum
costs $O(mb)$ edges. The contraction theorem then gives query work
$2^{O(mb/\sqrt{\log b})}$. Fresh mixed-target labels and a matrix-only
dependency recover the target. With base listing, failed queries, seed
restarts, root extraction, prime selection and linear algebra charged, the
total is $2^{O(b/m+mb/\sqrt{\log b})}$, and $m\approx(\log b)^{1/4}$ gives
the displayed exponent. The hashed bases are generic in the sense of
\S\ref{sec:pg-obligations}, so the route's non-generic content is entirely
in the contraction theorem.

\paragraph{What the programme checked.} Two in-session readings of the
decisive contraction proof, D.1 and the full D.4 recurrence, checked packet
shrink, degree invariants, interpolation counts, description and evaluation
costs, precision absorption, coefficient-ring valuations, Cauchy charts,
Hermite recovery and packet factorisation, and found no concrete gap. The
parent session read the preprint's multiplication construction, the
point-addition encodings, the hash coverage and the complete-cost analysis,
and found no transfer gap. The review itself says these are scoped proof
readings, not certification of either manuscript or of any implementation.
Implementation checks are all small:
\begin{itemize}
\item the native tensor crate passes 34 release correctness tests,
  including forced recursive and direct controls and a tiny complete
  relation pipeline;
\item an independent GMP reconstruction unit passes 72 fixed polynomial
  cases, 308 Cauchy charts, 8{,}832 node-error and 1{,}152 final-error checks
  and four negative controls, with byte-identical replay; it is not an
  extension-ring multiplication or contraction test;
\item a Rust port of the Go \texttt{ecdltensor} prototype passes all 19 Go
  tests with byte-identical output, and both demos recover the planted
  scalar (\art{research/ecdl_tensor_ic_rust_port_20261029/}); its counters
  are exponential and it does not implement the contraction algorithm.
\end{itemize}

\paragraph{What was not measured.} The speed study
(\art{research/ecdl_tensor_ic_speed_20261030/}) built an
\texttt{ic.tensor\_contraction\_v1} arm in ecbench and states that the
measurement has not run; it predicts $S$ growing about like $r$, since one
relation costs about $\abs{F_1}\abs{F_2}\approx2r$ additions in the
prototype. The counts-only protocol v2 for the tensor and hashed-native
panels (\art{research/hashed_ecdl_counts_20261031/}) superseded v1 before any
data, and no session has run under either. The finite-size review
(\art{research/hashed_ecdl_finite_size_20261030/}) states that the
preprint's literal constants force a finite-size fallback on every
practical graph, which does not refute an asymptotic theorem; it proposes
search strategies below the recursion's onset without measuring any. Its
ECCp-131 fixture uses a certified hash modulus
$s=3102808443644027376164572183694178153449$.

\paragraph{Census and elimination.} Two of the skeptic-reviewed leads, the
L09 admission census and the L06 width bound, bear on this route; their
numbers are in \S\ref{sec:tn-leads}. The elimination campaign of 1
November (\art{research/prime_ic_elimination_20261101/}) re-fetched both
PDFs at their pinned commits, matched the hashes, and kept the route at
rank~1 of its frozen register: the native finite fallback is rejected;
symbolic guards first pass at $\lambda=2^{480}$; semantic cuts expose
154--354 raw fanout edges but only 7--15 exact signal classes; natural ranks
beat support-matched sparse nulls, while all four matched cap-1 typed
topology nulls have exactly the same rank. Its state for the route is
``legacy census infeasible; minimal typed global-wiring mechanism
rejected; theorem/full null open''. A rank-census design
(\art{research/tensor_boolean_breakthrough_20261031/}) proposes to test
whether natural boundary tensors have lower flattening rank than rewired
nulls before any solver search; its E0 core is implemented and nothing
affected has been measured. The source-inspection audit of implicit tensor
entries (\art{research/hashed_ecdl_entry_cost_20261104/}) is claim
\texttt{PCC-20261104-1df881}, and the structured-predicate integration
report is \texttt{PCC-20261104-458dda}; both are of unverified standing.

\paragraph{What would test it.} The rank-census design of 31 October
stages the test as six gated experiments plus one SAT/BDD triangulation,
none authorised to spend benchmark budget until its predecessor's gate
passes:
\begin{itemize}
\item E0, canonical IR equivalence and mutation: every compared network
  regenerated, hashed and replayed with its semantics unchanged;
\item E1, a boundary-tensor flattening-rank census against graph- and
  predicate-preserving rewired nulls, with a held-out target and null gap;
\item E1S, SAT and BDD triangulation on the same pairs, to see whether exact
  Boolean solvers exploit the E1 signature;
\item E2, local Fourier and partition-rank certificates that predict at
  least 20\% lower held-out exact work;
\item E3, certified rewrites and orders with at least a $2\times$ held-out
  gain and a lower growth slope;
\item E4, a structure-specific contraction engine beating the strongest
  exact baseline by $2\times$ in operations and memory;
\item E5, an end-to-end relation oracle improving on four increasing sizes
  and an unseen target, still with no asymptotic claim.
\end{itemize}
The finite-size review proposes five ways to use the hashed relation
language below the recursion's onset: making inverse-hash selectors the free
inputs, solving for one witness instead of counting completions,
width-aware encoding and contraction plans, instance-certified recursion
with small packet parameters, and nested bases with partial relations. It
notes that enumerating every selector tuple remains generic decomposition
search. None of the seven experiments and none of the five routes has a
result.

\paragraph{Reading.} The route is the programme's best-supported
\emph{claimed} sub-exponential algorithm for prime-field ECDLP, and it is
external prior art. Its correctness rests on an external result that the
programme read but did not verify, its constants make it infeasible at
finite size as written, and it has no measurement. It does not move a row
of any table in this treatise, and the iterative success criterion does not
apply to it until a counted cost exists. The norm and divisor candidate of
\S\ref{sec:pa-divisor} does not inherit its compactness: its literal
translation has $O(mb^2)$ edges.

\section{EcGFp5 and EcMasFp5 over $\FF_{p^5}$}\label{sec:pt-gfp5}

EcGFp5 and EcMasFp5 are curves over $\FF_{p^5}$ with $p=2^{64}-2^{32}+1$
and $\FF_{p^5}=\FF_p[z]/(z^5-3)$. EcGFp5 is a double-odd curve with one
rational 2-torsion point. Chapter~\ref{ch:summer} records that the
September lane on EcGFp5 symmetrisation measured no $m=5$ ideal degree
(\texttt{DEC-20260923-d86ff2}). That statement is superseded: the 1
November solve below measured $D=2^{16}$ on EcGFp5-shaped curves over toy
primes. It is still true that nothing was measured at the 64-bit $p$.
Chapter~\ref{ch:extension} gives the Joux--Vitse $(n-1)$-decomposition at
$n=5$ over $E(\FF_{p^5})$ for toy $p$.

\paragraph{Desk evaluation.} The four atlas leads were evaluated on 31
October (\art{research/ecgfp5_lead_evaluation_20261031/}; desk evaluation,
every cost a model from cited constants). The design note's ``$\ge2^{142}$''
multiplies the large-prime system count $p^{2-2/5}=2^{102.4}$ by a
per-system cost without the $5!=120$ inverse decomposition probability, and
is not a lower bound for the Gaudry--Diem family. The cheapest member is
the basic algorithm with a factor base tuned to the per-system cost $c$,
whose total scales as $c^{1/3}$, about $2^{115.5+\log_2(c)/3}$ 64-bit word
operations: $2^{126}$--$2^{136}$ for the repository's per-system bands. In
raw time every variant is 32--45 bits below rho (rho $\approx2^{159.3}$
steps, $2^{169}$ word operations), but the tuned base has $2^{57}$--$2^{63}$
columns and needs $2^{64}$--$2^{71}$ bytes against $2^{35}$--$2^{45}$ for rho
with distinguished points; capped at a petabyte the method is about $2^{27}$
above rho. In a memory-aware model rho remains the better attack. The
verdicts are: symmetrised Gaudry \cls{revise} (measure $c$, not $D$);
Joux--Vitse $(n-1)$ for the full logarithm~\cite{jouxvitse2012}
\cls{close} ($(n-1)!\,\tilde c\,q^2/2=2^{131.6}\tilde c$ with
$\tilde c\approx2^{33.5}$--$2^{34.4}$ cycles lands within $\pm2.5$ bits of
rho, at about $2^{71}$ bytes, and loses to rho for any base of at most
$2^{62.5}$); the odd-characteristic cover
\cls{close} (genus 65 for EcGFp5 and 97 for EcMasFp5, index calculus at least
$2^{206}$ on the best constructed cover; the atlas keeps this lead
\texttt{open\_blocked} until the source-freeze block is lifted); oracle-assisted static
Diffie--Hellman \cls{revise} (about $2^{63}$ oracle queries, and no deployed
service found that exposes the oracle).

\paragraph{The quotient solve.} The decisive experiment
(\art{research/ecgfp5_quotient_solve_20261101/}; \art{benchmarks/20261101-ecgfp5-quotient-solve/};
pre-registered, amendments A1--A5) solved the $\FF_p$-rational
$(\ZZ/2)^4\rtimes S_5$ quotient of the five-summand system on EcGFp5-shaped
curves at $p'=4111$, $262231$ and $33554371$ with msolve 0.6.5, on one
4-vCPU host. It found $D=2^{16}=65{,}536$ at all three primes, exactly the
affine B\'ezout bound $8^5\cdot2$, confirming $D=2^{(m-1)^2}$ with no
reduction below B\'ezout; the planted decomposition was recovered at every
prime. One complete solve costs $c_{\rm cyc}=2^{46.30}$, $2^{46.56}$ and
$2^{46.63}$ cycles. The evaluator's frozen gate G1 reads \cls{close} when
$c_{\rm cyc}\ge2^{46}$, so all three cells close, by margins of $0.30$,
$0.56$ and $0.63$ bits. The F4 work is $p$-independent (the same final
$5578\times71114$ matrix at every prime), F4 alone is $2^{44.72}$--$2^{45.29}$,
and about two-thirds of $c$ is msolve's rational parametrisation. Weighted
orders and signature-based bases did not lower $c$. At $m=3$ all five arms
gave the predicted degrees in 30 solves, and msolve was found to crash for
every prime from 65537 to 253229, a solver defect, so 65551 was replaced
before any $m=5$ run.

\paragraph{The block-Wiedemann finish.} The finish was then replaced by a
block-Wiedemann solve~\cite{coppersmith1994} on the BNS22 chain matrix of the
reduced grevlex basis (\art{research/ecgfp5_quotient_bw_finish_20261102/};
\art{benchmarks/20261102-ecgfp5-bw-finish/}; frozen before the first cell,
no amendments). The eliminant equals msolve's coefficient for coefficient,
the stability check passes on 4{,}920 chains, and the planted point is
recovered at every prime. Whole-solve CPU falls $1.83$, $1.97$ and
$1.72\times$ at the three primes, clearing the frozen $1.5\times$ bar, so the
finish is the new per-system baseline; the finish alone is $3.15$--$4.12\times$
cheaper and the F4 work is unchanged. The new $c_{\rm cyc}$ is $2^{45.43}$,
$2^{45.58}$ and $2^{45.84}$, below the \cls{close} line by $0.57$, $0.42$ and
$0.16$ bits, but the frozen $[c/4,4c]$ band straddles the line at every
prime, so the verdict is unchanged (Table~\ref{tab:pt-gfp5}). The distance to the \cls{promote} line
of $2^{38}$ is $7.43$, $7.58$ and $7.84$ bits. This is an engineering
improvement of a per-system stage cost at toy primes, not an end-to-end
ECDLP cost. An earlier HNF-based finish was a memory-bound null at the
development prime $p=331$: the determinant step was killed at
13.33\,GiB on a 15.7\,GiB host
(\art{research/ecgfp5_quotient_hnf_finish_20261102/}).

\begin{table}[t]
\centering\small
\begin{tabular}{@{}rrrrrl@{}}
\toprule
$p'$ & W0 CPU-s & BW CPU-s & speed-up & $\log_2c_{\rm cyc}$ W0 / BW & G1 band (BW) \\
\midrule
4111 & 31{,}004 & 16{,}912 & $1.83\times$ & $46.30$ / $45.43$ & $[43.43,47.43]$ \\
262231 & 37{,}036 & 18{,}823 & $1.97\times$ & $46.56$ / $45.58$ & $[43.58,47.58]$ \\
33554371 & 38{,}778 & 22{,}516 & $1.72\times$ & $46.63$ / $45.84$ & $[43.84,47.84]$ \\
\bottomrule
\end{tabular}
\caption[]{One complete $m=5$ solve of the EcGFp5-shaped quotient system at
three toy primes, $D=2^{16}$ in every cell: msolve with rational
parametrisation (W0) against the block-Wiedemann finish (BW); one 4-vCPU
host, $c_{\rm cyc}=\text{CPU-s}\times2.8\cdot10^9$. The \cls{close} line is
$2^{46}$ and the \cls{promote} line $2^{38}$. Sources:
\art{research/ecgfp5_quotient_solve_20261101/},
\art{research/ecgfp5_quotient_bw_finish_20261102/}.}
\label{tab:pt-gfp5}
\end{table}

\paragraph{Representations.} A desk audit of 2 November
(\art{research/ecgfp5_representation_audit_20261102/}) closed every
degree-2 coordinate change and every larger symmetry group: EcGFp5 has one
rational 2-torsion point and no rational 4-torsion, so the only usable group
gives $D=2^{16}$ at $m=5$. It found one open lever, Pornin's $u=x/y$: since
$u(P+N)=-u(P)$ and $u(-P)=-u(P)$, the quotient is a polynomial-ring form in
$s_k=e_k(u_i^2)$ and $\pi=\prod u_i$, with measured $D=2^{(m-1)^2}$ at $m=3$
and $m=4$ (16 and 512) and the planted point recovered. That corrected a
ledger statement: \texttt{CORR-20261102-7e3a91} scopes the earlier claim that
the ring form is ``unavailable'' for double-odd curves to coordinates with
$\varphi(-P)=\varphi(P)$. Any gain from $u$ is expected to be a constant
factor. The atlas lead \texttt{LEAD-ICP-gfp5-symmetrized-gaudry} is
\texttt{closed\_dominated}, with the block-Wiedemann
finish as a measured per-system improvement still short of the
\cls{promote} bar.

\paragraph{$\FF_{p^4}$.} Two items extend \S\ref{sec:gaudry}. The $D_4$
index calculus ran end to end in ecbench on $E(\FF_{29^4})$,
$\#E=708{,}208=16\cdot44{,}263$, with every phase charged: $S=7.17\cdot10^8$
against $1.065$ for negation rho, audit OK on three of three records
(\art{research/gaudry_quartic_d4_ecbench_20261013/}). A later leaderboard
engineering row cut the plain-$\Sem5$ cost $C_4$ from $1.209\cdot10^{12}$ to
$8.69\cdot10^{10}$ $\FF_p$ multiplications ($13.90\times$) with sparsest
pivots, chain eigensteps and shape replay
(\art{benchmarks/20261019-s5-plain-c4-options/}), where shape replay alone
had been a 2--5\% lever (\art{research/gaudry_quartic_replay_20261013/}).

\section{ECCp-131}\label{sec:pt-eccp131}

Chapter~\ref{sec:prime} gives the ECCp-131 parameters, the zero-relation
probe and the rho rates. Four further studies are nulls or constructions.
\begin{itemize}
\item \emph{A full-width auxiliary tower.} The Kummer tower cannot have
  $t\ge2$ on the ECCp-131 field, since $v_2(p-1)=1$. A rational eight-step
  2-isogeny tower was built instead on the auxiliary curve
  $E':y^2=x^3+x+301$ over the exact field
  (\art{research/eccp131_full_width_20261023/REPORT.md}; of 301 auxiliary
  orders computed, two were divisible by 256, at $b=278$ and $b=301$, and the
  curve with $b=301$ was accepted): 256
  distinct coset abscissae, 8 of 8 V\'elu maps checked, 115 abscissae
  lifting to the published curve, 26{,}450 signed pair points, 8 of 8
  planted controls, and zero natural setup or target rows, rank zero.
\item \emph{A multiplicative-coset screen.} Twelve $g\mu_5$ cosets and,
  in a follow-up, sixteen $g\mu_{10}$ cosets ($\mu_{10}$ is the maximal
  smooth subgroup, $p-1=2\cdot5\cdot r$ with $r$ prime) against equal-size
  small-abscissa bases found zero natural rows and rank zero for both seeds,
  with every planted control recovered
  (\art{research/eccp131_structured_ic_20261022/REPORT.md}).
\item \emph{A representation filter.} A Wagner-style coordinate-bit filter
  trained on true four-sum relations against a rotated null, with 3{,}450
  configurations in each of six base/size cells
  (\art{research/eccp131_representation_filter_20261102/REPORT.md};
  \art{benchmarks/20261102-eccp131-representation-filter/}), qualified in
  no held-out cell: every recorded lower confidence ratio is below 1, and
  the two superficially positive small-$x$ winners at $B=64$ and $B=256$ are
  post-addition filters that need all $\Theta(B^2)$ pair sums. The verdict is
  \texttt{no\_qualifying\_filter\_in\_frozen\_family}, with 0 rows and rank 0
  at arity 2 and 3.
\item \emph{A CUDA port.} The ECCp-131 and ECC2-131 rho walks were written
  as per-thread functions that compile as device code
  (\art{research/certicom131_cuda_port_20261023/}). The kernels compile and
  assemble for sm\_80--sm\_120 but have never executed on a GPU; there is no
  device parity result and no GPU rate.
\end{itemize}
The full-parameter index calculus remains the open problem
\art{problems/eccp131-full-parameter-index-calculus.md}.

\section{Audits of standard prime curves}\label{sec:pt-standard}

Chapter~\ref{sec:prime} lists the audits of the NIST, SECG and CryptoPro
curves that found no ECDLP weakness. The later ones concern curve
provenance and class-group statistics, and none bears on ECDLP hardness.
\begin{itemize}
\item \emph{The NIST seed skip rule.} Of the 21 seed candidates the 1997--98
  runs accepted under every known criterion, 10 were published. None of 21
  new preregistered binary features explains the 11 skips, after about 60
  earlier features were ruled out
  (\art{research/nist_seed_skip_filter_20261017/}). A later search found an
  independently checked two-test empirical separator for all 21 eligible
  curves, with a root-convention dependence and a calibration that blocks
  reading it as the historical explanation
  (\art{research/nist_skip_criterion_20261025/}). The same question for the
  binary B- and K-curves is a measured census
  (\art{research/nist_binary_skipped_20261020/}), and a factorisation
  continuation fully factored eight of the nine distinct integers left by
  an external DiSSECT run, leaving one 444-bit composite cofactor of a
  P-256 extension-twist order (\art{research/nist_skipped_curve_dissection_20261021/}).
\item \emph{P-224's class group.} The 2-rank of $\mathrm{Cl}(D_K)$ is
  exactly 6 for P-224 (4 for P-256), and the class number is estimated at
  $3.5185\cdot10^{32}=2^{108.1167}$, with a GRH-conditional interval
  $[2^{108.098},2^{108.135}]$ under an assumed constant; the exact value is
  not known (\art{research/p224_class_number_20261022/}). A 2-rank of 6 is
  not unusual: P-224 sits at the 91.5th mid-rank percentile of random
  224-bit primes, upper tail 12.3\% $[10.3,14.5]$
  (\art{research/p224_two_rank_null_20261022/}).
\item \emph{Cheon's attack on P-224's shape.} The DLP with auxiliary inputs
  was scaled on P-224-shaped curves
  (\art{research/cheon_p224_shape_20261022/REPORT.md}). The pre-registered
  rho-reliability gate H3 failed, with 5 of 84 sealed rho runs outside
  $S\in[0.3,4]$, so every Cheon/rho ratio and crossing in that report is
  withdrawn; the Cheon-only scaling stands. It is in any case not plain
  ECDLP.
\item \emph{Special prime shapes.} Curves built by the CM method over
  192-bit Proth, cyclotomic $\Phi_d(m)$ and smooth $p\pm1$ primes all have
  embedding degree $k>4000$, like the random control; only pairing-friendly
  constructions give $k$ of 10--24 (\art{research/special_prime_curves_20261023/}).
\item \emph{Rho on NIST curves.} Stepping to the better of $P+M_j$ and
  $P-M_j$ in the \texttt{gpu/nist} small-batch kernel stays disabled after an
  audit; the estimates are historical diagnostics, not measured solve costs
  (\art{research/nist_best2_small_batch_20261027/}).
\end{itemize}

\section{What the three chapters settle}\label{sec:pt-summary}

On every prime curve measured, rho with negation is cheaper than every
index-calculus arm, generic or algebraic. The generic track's best arms are
the residual collectors at $1.3$--$1.7\times$ rho ($w=0$; $1.54$--$1.64$ in
the cells with a validator report), a constant at rho's exponent with no
growth over 36--42 bits; the $k$-sum tables are
$5.4$--$11.0\times$ rho and above baby-step giant-step. The algebraic track's
priced prime rows are the Kummer-tower rows and the Curve25519 F4 arms; the
best of the latter, with a validator report, is still $1.83\cdot10^4\times$
rho at $w=0.1$, and the cost scales as $r^{1.03}$ (\S\ref{sec:c25519}). Its degree
studies show growth or input-pinned degrees in every presentation, with
matrices that grow when degrees do not. The one claimed sub-exponential
route is external, conditional on an unverified theorem, infeasible at
finite size as written and unmeasured. On EcGFp5 the measured per-system
cost moved by $1.72$--$1.97\times$ and remains about 7.4--7.8 bits from the
promote line.

The chapters state these corrections and withdrawals:
\begin{itemize}
\item the residual collectors' $w>0$ gains were the price of a byte
  counter, withdrawn under option~B (1 November); the scale study's C1 tally
  falls from 6 of 8 to 2 of 8, and a 9-versus-8 tally of 12 cells is
  unreconciled (\S\ref{sec:pg-bytes});
\item the memo gate ``$e<1$ at $m=4$'' is replaced by $s<m-3$
  (\S\ref{sec:pg-gate});
\item the cofactor-cleared prime base is \texttt{invalid\_controls}, not a
  win (\S\ref{sec:pg-cofactor});
\item $N=64$ at 20 bits is unresolved, and the 24- and 28-bit base
  selections are relabellings (Chapter~\ref{ch:distance});
\item the small-root margin of $-0.048$ is superseded by about $-0.6$
  (\texttt{CORR-20261024-b2a799}, \S\ref{sec:pa-smallroot});
\item the September statement that no $m=5$ EcGFp5 degree was measured is
  superseded by $D=2^{16}$ at toy primes, and the ring-form statement is
  scoped by \texttt{CORR-20261102-7e3a91} (\S\ref{sec:pt-gfp5});
\item the symmetric-tower rows are instrument diagnostics, not natural
  prime-order benchmarks (\S\ref{sec:pa-pkm});
\item the Cheon/rho ratios on P-224's shape are withdrawn
  (\S\ref{sec:pt-standard}).
\end{itemize}
The open items are the discriminating Kummer cells ($m=3$, $N=21$
symmetric; $t=14$ at cap 6), the algebraic reference arm, a counted cost
for the tensor route, the $\FF_p$ CDCL(T) propagator, the APS complete-cost
panel, and a proved uniform bound for a divisor-norm oracle.
